% SOURCE docs/release/2026-10-29/source/docs/writing-repair/2026-10-29/src/f60c8396c1dea5f9.tex CHAR 0
\chapter{Understanding the records behind an estimate}
\label{exec16:datahub}

A bank considering a new servicing policy first needs to locate the histories
that describe eligibility, actions, cash flows and observation. An enterprise
catalog can identify datasets, owners, schema fields and transformations that
would otherwise have to be rediscovered. The causal proposal should reuse that
work. It must also preserve the distinction between a record about a dataset
and evidence about the decision made for a particular customer. The chapter follows that distinction through catalogues, table preparation, historical features, computation and lineage, then validation and uncertain linkage. DataHub and OpenMetadata address catalogue context; DuckDB and Polars address transformations; Feast addresses feature retrieval; OpenLineage and the workflow comparison address execution history. Pandera and Splink show why a well-formed table can still represent the wrong population or an uncertain customer map. These components perform complementary tasks and require different acceptance evidence.

DataHub is a concrete candidate for this boundary. Its open repository contains
typed dataset and field identities, ownership and lineage aspects, a Python
client and metadata ingestion tooling. The inspected revision, dated 7 October
2026, pins sixteen selected original files. The Python distribution is named
\texttt{acryl-datahub}; its version is dynamically supplied, so the source
revision does not establish the identity of an installed distribution. The
README separately describes OSS and Cloud. The referenced MCP server and
Cloud SDK require their own package, rights and behavior review. No service,
SDK, connector or MCP server was executed for the comparison below.
\citep{exec16_dh_primary,exec16_dh_rights}

\section{The metadata graph has a useful, specific meaning}

The dataset constructor takes a platform, dataset name, platform instance and
environment, then delegates the construction of its URN. Field identity adds
a field path to a dataset identity. An exporter should preserve those opaque
identities and their versions. The display name \texttt{balance} is insufficient:
three stipulated fields in production/east, development/east and production/west
are distinct even when their dataset and leaf field names agree. An exact
identity example has three complete tuples but only one leaf name. Merging
them changes which population and source produced an observation; a downstream
estimator cannot reconstruct the lost distinction.
\citep{exec16_dh_primary,exec16_dh_exact}

The \texttt{UpstreamLineage} aspect represents upstream datasets and optional
fine-grained field relationships. \texttt{FineGrainedLineage} can retain
transformation text and a query reference. These fields help explain how a
derived value was computed. The type also has a \texttt{confidenceScore}
whose default is one. The inspected definition calls it confidence in lineage;
it supplies no calibrated probability that a corresponding causal arrow is
true. Its optional matching status describes the resolution of a URN against
the catalog at ingestion time, and the comments explicitly say that this
status is not re-evaluated later. A later exporter must preserve the distinction
between exact, normalized, unresolved and unassessed references.
\citep{exec16_dh_primary}

This distinction matters when an agent helps construct a diagram. A pipeline
may transform a contact table and a recovery table into an analysis table.
That lineage supports a provenance statement. It does not establish that
contact caused recovery, that contact assignment was unconfounded, or that
every predictor was available before contact. A description or glossary term
can suggest a variable interpretation, but the proposal should retain its
source, owner and unresolved alternatives until the causal interpretation has
separate support. The agent may propose a mapping; it must not promote a
lineage-confidence default into scientific certainty.
\citep{exec16_dh_primary,exec16_dh_contract}

\section{Identical catalog context can support different interventions}

An exact generated example makes the missing information visible. Let $A$ be
a binary offer and $Y$ a binary recovery. In a first world, $A$ is a fair
random draw and
\[
 Y=A\mathbin{\oplus}E,\qquad P(E=1)=1/4,
\]
where $\oplus$ denotes exclusive-or and $E$ is independent of assignment.
In a second world, an unobserved fair borrower state $U$ determines assignment,
$A=U$, and $Y=U\mathbin{\oplus}E$; changing the offer leaves the recovery
mechanism unchanged. Both worlds give the strictly positive observed table
\[
 \begin{array}{c|cc}
 &Y=0&Y=1\\\hline
 A=0&3/8&1/8\\
 A=1&1/8&3/8
 \end{array}
\]
and can use the same two source tables, field descriptions and transformation
lineage. Direct enumeration gives
\[
 \begin{aligned}
 E[Y\mid\operatorname{do}(A=1)]-E[Y\mid\operatorname{do}(A=0)]
 &=1/2 &&\text{in the first world},\\
 &=0 &&\text{in the second world}.
 \end{aligned}
\]
The observed conditional difference is $1/2$ in both. At a stipulated USD~100
per recovery and USD~30 per offer, the incremental net values are USD~20 and
USD~$-30$. The opposing decisions arise with identical observed tables and
catalog structure. This is an authored information example, not a bank effect
estimate or a failure observed in DataHub. More catalog entries cannot alone
choose the assignment mechanism. \citep{exec16_dh_exact}

Useful catalog integration therefore retains the evidence for an assignment
claim alongside the variable mapping. For a randomized campaign, that may
eventually include the actual allocation record and deviations. For an
observational proposal, it may be a declared adjustment argument with competing
graphs and sensitivity analyses. An empty assignment-evidence field remains
unknown. It cannot be filled by a well-connected asset graph.
\citep{exec16_dh_contract}

\section{Separate the clocks before constructing a historical feature}

The inspected audit-stamp type records when a resource or association entered
a lifecycle stage and who acted. An upstream lineage stamp records who
reported that lineage and when. Ownership has an audit stamp whose zero time
explicitly denotes missing data. Schema-field creation and modification stamps
likewise describe metadata objects. None of these definitions, by itself,
states when a particular customer-level value first became available to a
servicing decision. \citep{exec16_dh_primary}

Let $t$ be the action time, $e_j$ a feature's event time, $a_j$ the time its
value became available, and $m_j$ the time the catalog observed its metadata.
A proposed historical input contract requires $a_j\leq t$ for a feature used
by the policy, together with evidence for the relevant version. The inequality
$e_j\leq t$ is insufficient when an event is recorded or revised later.
Similarly, $m_j>t$ need not invalidate a retrospectively acquired description
of an older value; it means that the exporter needs an additional historical
availability argument. The clocks have different meanings.
\citep{exec16_dh_contract}

For a generated example, take an independent fair baseline $B$ and recovery
$Y$ at day 30. The baseline is available on day $-1$. A derived field $Z=Y$
is labeled with an origination event on day zero but becomes available only
on day 31. Both fields have schema creation time $-60$ and their catalog
snapshot is acquired on day 40. Enumerating the four equally likely $(B,Y)$
pairs gives classification agreement $P(B=Y)=1/2$ and $P(Z=Y)=1$.
The latter is perfect retrospective leakage. At a day-zero action, the
availability rule admits only $B$; filtering solely by event time also admits
$Z$. A third field with unknown availability stays unresolved even if its
event time is old. These exact agreements describe stipulated variables,
not a fitted prediction benchmark. \citep{exec16_dh_exact}

DataHub's inspected \texttt{get\_aspect} method supports an explicit version,
with version zero documented as latest. A latest schema can be useful for
discovery, but it is not automatically the schema used by a historical policy.
The proposed export should bind the requested aspect version, returned bytes,
source-data version and clock evidence. A sequence of reads is not declared
an atomic catalog snapshot merely because every request succeeds. Version
consistency and concurrent changes belong in the prototype acceptance cases.
\citep{exec16_dh_primary,exec16_dh_contract}

\section{A bounded export must preserve incompleteness}

The selected aspect reader distinguishes a missing aspect, which yields
\texttt{None} for HTTP 404, from other HTTP failures and from a malformed
response lacking the expected aspect. It also rejects a time-series aspect
on this particular path. An adapter should retain these distinctions. Missing
metadata is not evidence that an asset has no owner, that a covariate has no
confounders, or that the underlying data were complete. An authorization or
transport failure must not become an empty successful causal context.
\citep{exec16_dh_primary}

Lineage traversal has two further source-level details. The inspected client
maps requested hop counts above two to the categories one, two and ``3+''.
That is not a strict three-edge radius: a generated chain with distances
one through four retains distance four under the ``3+'' category, whereas a
literal three-hop contract would exclude it. The method's \texttt{count}
argument is passed to a scrolling query, and the implementation accumulates
results while a continuation cursor is returned. Two generated pages of 500
and one records therefore produce 501 accumulated records. The per-query
count is not a proved whole-export bound. These are static interpretations
of the saved code, not observations from a running server.
\citep{exec16_dh_primary,exec16_dh_exact}

A proposed export wrapper should make its own request, result, byte and elapsed
time limits explicit. It should track cursors, preserve each returned page,
and mark exhaustion, repetition or inconsistency as partial or failed retrieval.
In the authored cursor example, seeing the same nonempty cursor twice produces
a repeated-cursor disposition after two requests, with completeness false.
No live failure is alleged. The example specifies what a bounded wrapper must
do if that response sequence occurs. A partial graph cannot support the
claim that all relevant upstream sources have been found.
\citep{exec16_dh_exact,exec16_dh_contract}

Lifecycle and projection also affect completeness. The inspected status model
separates the legacy removal flag from a lifecycle stage whose search visibility
may differ; its comments warn that these are not automatically synchronized.
The lineage-result conversion retains schema-field path entries while omitting
other path-node types, such as queries. A causal provenance export should
therefore retain raw aspect/query evidence and lifecycle context, in addition
to a convenient summarized path. It should not silently infer a missing
transformation from the projected result. \citep{exec16_dh_primary}

\section{A reuse choice with a concrete prototype}

The present choice is to retain DataHub as a pinned technical reference and
a candidate for conditional integration with an existing catalog. Adopting
the entire platform as a required causal engine would add a deployment and
maintenance dependency before the narrow use has been justified. Porting its
catalog or client into the causal core would create another implementation
to keep compatible. A small export boundary, with a manually prepared public
metadata fixture as the initial comparator, is a more testable proposal.
Other catalog families remain separate candidates; their names do not supply
equivalent APIs, rights or clock semantics. \citep{exec16_dh_contract}

The root source license is Apache 2.0 and the repository has a substantial
NOTICE file. The inspected redistribution conditions address license copies,
modification notices, attribution and applicable NOTICE content. The Python
project declares Python 3.10 or newer and numerous dependencies, including
version ranges; its dependency file says it is generated from another source
of truth. This review does not establish a complete resolved distribution,
dependency-license inventory or equivalence to that generating file. The
source revision and a recent commit show observed maintenance activity, not
a support commitment or compatibility with a future bank's server.
\citep{exec16_dh_rights}

The prototype should first consume a frozen, project-authored export containing
opaque entity and field identities, aspect versions, raw hashes, acquisition
time, lineage evidence, lifecycle state, unknown fields and separate feature
availability declarations. It should reproduce the four exact example families,
retain the five explicit missing-context cases, and refuse to emit causal arrows
solely from catalog lineage. A later source-matched SDK fixture must check
actual serialization, pagination, schema changes, deletion visibility and
failure behavior before the proposed adapter receives integration permission.
No such permission is granted by the present reference decision.
\citep{exec16_dh_contract}

An agent using the library, CLI or MCP interface should receive the same
bounded context object and evidence links. Read operations and proposed metadata
changes should have separate interfaces: the inspected SDK contains both query
and mutation methods, including automatic column matching. A discovery request
should not invoke a lineage mutation as a side effect of trying to understand
the data. Descriptions remain source material, not instructions to change the
analysis or credentials. This is a proposed consumer contract; it does not
claim that an existing MCP integration implements it.
\citep{exec16_dh_primary,exec16_dh_contract}

The catalog owner would eventually supply the permitted scope, access route
and relevant historical evidence. The prototype maintainer owns source/version
matching, conversion and bounded failure handling; a scientific reviewer owns
the assignment and measurement claims. The economic comparison includes
manual mapping time, maintenance, infrastructure, entitlement, semantic review
and incident investigation. A public fixture can establish whether the proposed
boundary prevents identifiable errors before any bank connection is sought.
It cannot establish actual owner truth, causal identification or bank return
on investment. \citep{exec16_dh_contract}
\label{exec16:datahub-end}


% SOURCE docs/release/2026-10-29/source/docs/writing-repair/2026-10-29/src/919c01d751f6bd4a.tex CHAR 0
\section{Analytical preparation must preserve the population and its information clock}
\label{exec19:duckdb-preparation}

A servicing analyst may ask whether an additional contact improves recovery,
then construct the analysis table by joining account records, contact events
and repayment observations. The resulting query can execute successfully while
changing the population, discarding unobserved outcomes or incorporating
information that did not exist when contact was chosen. The database is
therefore part of the causal workflow's measurement boundary. Its useful
product role is to make a declared preparation repeatable and inspectable;
query success cannot establish the causal premises of the subsequent analysis.

DuckDB is a concrete candidate for local analytical preparation. The selected
historical source is version 1.4.0, released on 16 September 2025, at commit
\texttt{b8a06e4a2267}. This choice fixes a reproducible object for discussion;
it is not a claim about the newest release, present support or any installed
binary. The inspected C++ interface distinguishes a materialized
\texttt{Query} call from streaming and pending paths. Result objects expose
column information, errors and potentially chained results, while fetching a
materialized result consumes its contents. These interfaces suggest a narrow
first integration: a typed request produces one bounded, completed table and
an explicit completion status. They do not justify treating every query handle,
first batch or returned row as a completed analysis. The Python client is a
separate wrapper boundary whose conversions and installed dependencies require
their own inspection.~\citep{exec19_duckdb_source}

\subsection{A missing outcome is not an observed failure}
\label{exec19:duckdb-missingness}

Consider an authored six-account example with three contacted and three
uncontacted accounts. The contacted outcomes are $(1,\mathrm{missing},
\mathrm{missing})$ and the uncontacted outcomes are $(1,0,0)$, where one means
recovery. These are stipulated binary observations, without a randomization,
exchangeability or missingness assumption. The contrast among observed
outcomes is
\begin{equation}
 \widehat d_{\mathrm{observed}}=1-\frac13=\frac23.
 \label{exec19:duckdb-observed}
\end{equation}
If the two missing contacted outcomes are denoted $u,v\in\{0,1\}$, the
full six-account descriptive contrast instead satisfies
\begin{equation}
 d(u,v)=\frac{1+u+v}{3}-\frac13
       =\frac{u+v}{3}\in\left\{0,\frac13,\frac23\right\}.
 \label{exec19:duckdb-completions}
\end{equation}
Both endpoints are attainable, so $[0,2/3]$ is the sharp endpoint range for
this stipulated finite completion problem. The feasible set has three values;
it is not every real number in that interval. Neither the complete-case
contrast nor the completion set is a causal effect without additional
assignment and observation assumptions. Exact rational enumeration reproduces
all four binary completions.~\citep{exec19_duckdb_reference}

The economic consequence of the transformation is visible before any bank
experiment. Suppose solely for illustration that a unit change in this
descriptive contrast is valued at USD 300 and an additional contact policy
costs USD 40 per case. The resulting arithmetic is
\begin{equation}
 V(u,v)=300d(u,v)-40\in\{-40,60,160\}\quad\text{USD per case}.
 \label{exec19:duckdb-value}
\end{equation}
Using only observed outcomes produces the optimistic plug-in value USD 160;
the declared completion problem also permits minus USD 40. This is an
illustration of how an unacknowledged denominator choice can reverse an
economic recommendation, not evidence of treatment value or bank ROI. A
missingness model, additional observations or a deliberately restricted target
population would change the problem and must be stated before interpretation.
\citep{exec19_duckdb_reference}

The selected upstream aggregate tests reinforce why the interface must retain
these distinctions. On an empty input, the tests expect zero counts but NULL
for sums, averages and several other aggregates; grouped all-null values also
produce NULL aggregates. Those are inspected upstream expectations, not native
test results obtained here. A zero-row relation, two rows with no observed
outcome, and two observed zeros therefore require different records. The
proposed record carries total rows, observed outcomes and nullable summaries;
it never silently replaces a missing mean with zero. NaN, overflow and invalid
type conversions need additional explicit cases rather than being folded into
the SQL NULL convention.~\citep{exec19_duckdb_source,exec19_duckdb_reference}

\subsection{Joins and clocks define the observation being analyzed}
\label{exec19:duckdb-keys}

For a second exact example, let two contacted accounts have observed amounts
USD 100 and USD 0, and two uncontacted accounts both have USD 0. The
account-level descriptive difference is USD 50. Suppose a one-to-many contact
join duplicates the positive contacted account three times while leaving each
other account once. The joined table has six rows and a contacted mean of
USD 75. No account outcome changed: the join changed the weights. The input
and output keys, expected join cardinality and analytical unit must therefore
be part of the preparation request. A deliberate contact-event target is
possible, but it is a different target and does not turn six joined rows into
six independent accounts.~\citep{exec19_duckdb_reference}

Time ordering creates a separate problem. In the authored reference, an
account has a feature with value zero on day one, an action on day two and a
feature with value one on day three. Selecting the latest unrestricted record
returns the postaction value; selecting the latest strictly preaction record
returns zero. A production definition must additionally distinguish event time
from availability time: a backdated event received after the decision was not
necessarily available to the decision maker. The proposed table carries both
clocks, the action cutoff, time-zone interpretation and an explicit tie-break
for events with equal timestamps. This clock discipline is an authored
requirement; the small reference does not validate a bank's record history.

Row position is an unsuitable replacement for those keys. The selected tests
make null ordering configurable, and an insertion-order test explicitly
disclaims which tuples an unordered limit selects when insertion-order
preservation is disabled. A reproducible export consequently specifies a total
order with unique tie-breaks and a declared NULL convention. All 24 permutations
of the four-account reference return the same rows after sorting by its unique
account key. That exact property belongs to the reference construction; native
parallel execution, collations and timestamp conversions remain future
conformance cases.~\citep{exec19_duckdb_source,exec19_duckdb_reference}

\subsection{Local execution needs explicit permissions and completion}
\label{exec19:duckdb-controls}

The selected source initially enables external access. When external access is
disabled on a running instance, its setting handler adds already attached
database and write-ahead-log paths to allowed paths and permits a configured
temporary directory. Selected storage and attachment tests continue to permit
writes or checkpoints to attached databases after this setting is disabled.
It is therefore neither a read-only guarantee nor an operating-system sandbox.
The proposed worker needs independently defined filesystem permissions,
reviewed attachments and scratch storage, even when that database option is
false. Its input copies must be immutable under the declared execution policy.
\citep{exec19_duckdb_source}

Extension behavior is also a build and runtime question. The selected CMake
options default automatic installation and loading to false, but compile-time
definitions can change the configuration defaults. The base extension
configuration includes core functions and Parquet, and conditionally includes
jemalloc under its platform rules. Inspecting those declarations does not prove
the configuration of a wheel, executable or embedded library. For the initial
proposal profile, automatic extension installation and loading should be
disabled and any required artifact explicitly selected and bound. An agent
may propose a typed preparation request; it should not acquire unrestricted SQL
and installation authority as a side effect. The inspected autoload helper
catches exceptions and can return false. Such a dependency failure is not
evidence that an intervention is unidentified or that an estimated effect is
zero.~\citep{exec19_duckdb_source}

Transactions do not remove the need for those boundaries. The C++ begin,
commit and rollback wrappers inspect errors and throw them; a selected upstream
test expects a conflict for overlapping updates. These observations support
explicit transaction-error handling but do not prove a complete isolation
regime or an immutable snapshot of files outside the database. The preparation
manifest binds input bytes separately from the transaction identifier and
records whether an output is complete. A timeout, cancellation, incomplete
stream or later error invalidates a claim of completed output. Materialization
simplifies the initial completion boundary while increasing memory exposure;
row and byte limits, disk use and cancellation costs must be measured on the
selected prototype rather than assumed from the interface name.
\citep{exec19_duckdb_source}

\subsection{Reuse the original engine under a bounded proposal contract}
\label{exec19:duckdb-route}

The preferred proposal route retains the original engine behind that typed
preparation boundary. A port would reproduce SQL, optimization, storage and
conversion obligations without a demonstrated platform need. The selected
historical build declares C++11, twelve direct non-amalgamation third-party
subdirectories, additional conditional test directories and extension inputs.
The root MIT notice is not the complete distribution inventory. The inspected
RE2 and Zstd notices are BSD-style, utf8proc carries MIT and Unicode data
notices, and MbedTLS declares Apache-2.0 or GPL-2.0-or-later alternatives. These
four inspected third-party notices illustrate scope differences; they do not
clear every dependency, wrapper, extension or actual distribution arrangement.
\citep{exec19_duckdb_source}

A transparent explicit-file implementation remains the small-case reference.
SQLite is a separate SQL/storage candidate; dataframe and columnar libraries
are separate preparation candidates. Each alternative requires its own type,
join, missingness, ordering and resource review. The choice among them should
follow the declared workload and measured end-to-end costs, including frozen
inputs, conversions, failed requests, maintenance and retained evidence. No
unmeasured speed, feature parity or maintenance advantage is assumed here.

The proposal choice is consequently concrete while executable admission is
deferred. Before admitting the selected engine/client path, a prototype must
bind its complete distribution and notices, exercise native empty/all-missing/
zero cases, duplicate joins, information cutoffs, total ordering and precision,
and compare errors, cancellations, external-file changes and extension
behavior with the typed contract. The analyst owns the scientific population
and causal assumptions; the adapter maintainer owns conversion, execution and
resource behavior; the distribution owner reviews the actual retained software
and notices. None of these public or synthetic prototype steps requires bank
data. Institution-specific assignment, audit truth and economic value remain
distinct later questions. The database prepares the evidence; it does not
certify the causal meaning of that evidence.


% SOURCE docs/release/2026-10-29/source/docs/writing-repair/2026-10-29/src/35b2a95dc6edd027.tex CHAR 0
\subsection{Feast: reconstructing the information available to a banking decision}
\label{exec19:feast-clock}

Consider a credit team replaying yesterday's offer decisions. A transaction
can have an early business date, a later correction date and a still later
arrival in the system used by the offer engine. Joining on the business date
alone does not reproduce the information on which the original offer could
depend. Feast is a concrete candidate for organizing feature definitions and
retrieving historical and online features, but its usefulness for this
proposal depends on a precise observation contract. The selected original
revision is commit \texttt{9d42729d471a} of 7 October 2026. Its historical API
already provides an optional creation-time cutoff; the appropriate comparison
must acknowledge that capability and investigate its meaning for each
provider.~\citep{exec19_feast_source}

For a feature version $j$ associated with customer $i$, let $e_{ij}$ be its
event time, $c_{ij}$ its source creation time and $a_{ij}$ the first time that
the version, with the required transformations, was usable by the decision
system. Let $t_i$ denote the decision time and $\ell>0$ a declared feature
lifetime. The proposed admissible set is
\begin{equation}
 \mathcal J_i(t_i,\ell)=
 \{j:t_i-\ell\leq e_{ij}\leq t_i,\quad a_{ij}\leq t_i\}.
 \label{exec19:feast-eligible}
\end{equation}
All times use a declared timezone and ordering convention. If feature arrival
and action share a coarse timestamp, an additional sequence number or
conservative ordering is needed to establish which occurred first. The
inclusive inequalities in \eqref{exec19:feast-eligible} are meaningful only
after that convention is fixed. The latest admissible event and revision may
then be selected, provided that the revision order is well defined. A source
row with an unknown availability time does not establish membership in this
set. Conflicting values with the same event and revision coordinates require
an explicit conflict result or a separately justified source ordering.

The distinction is statistical as well as operational. If $\mathcal F_{t_i}$
is the information available to the decision maker, a deployable rule must
be measurable with respect to $\mathcal F_{t_i}$, with any independent policy
randomization stated separately. A retrospectively corrected covariate can
belong to a larger information set. Its association with a later outcome may
be useful for another scientific question, but it cannot silently replace a
predecision predictor, confounder or effect modifier. Nor does successful
retrieval establish exchangeability, positivity, a causal graph or calibrated
uncertainty. Those obligations remain with the analysis that consumes the
retrieved table.

\paragraph{What the selected API and providers actually specify.}
The historical API returns a retrieval job and defaults
\texttt{filter\_by\_created\_timestamp} to false. With the flag enabled,
the requested additional condition is $c_{ij}\leq t_i$ for views that declare
a creation-timestamp column; views without that column are unaffected. The
selected passthrough provider checks the offline store's declared support,
and the base store raises an unsupported-operation error when support is
absent. It would therefore be inaccurate to say that Feast necessarily
ignores creation time. It would also be inaccurate to equate that field with
$a_{ij}$ without evidence about ingestion delays, transformations and the
actual serving path.~\citep{exec19_feast_source}

The selected Dask implementation applies the event-time and lifetime filters,
then the optional creation cutoff, then duplicate handling. For late or
unknown creation times it blanks feature-side fields while retaining the
entity-side fields. The accompanying source tests specify the difference
between a later-created value of $0.6$ and an earlier value of $0.4$, exclusion
of null creation times, retention of an entity whose candidates are all late,
and separation of different entities. These are informative upstream test
specifications; they are not native test results obtained by this proposal.
The implementation treats a zero or unspecified lifetime as absence of a
finite lower event-time bound. The proposal must therefore distinguish an
explicitly unbounded history from a positive duration, rather than converting
zero into an instantaneous window.~\citep{exec19_feast_source}

The selected BigQuery template applies event-time conditions and, when
requested and configured, the creation cutoff before choosing the latest
event and creation timestamps. Its final join back to feature records uses
those timestamps and the entity-row identifier. If different records share
the chosen coordinates, the proposal cannot infer unique output merely from
the earlier ranking expression. A source uniqueness condition or a qualified
revision key is necessary. The Dask path also needs native checks for equal
timestamps and duplicate decision rows, since sorting and duplicate removal
do not by themselves establish the required population contract. This is a
specific conformance question, not a claim that a defect has been reproduced
in an installed provider.~\citep{exec19_feast_source}

\paragraph{An exact decision example with two kinds of late information.}
Consider the following authored finite law, chosen to make the information
constraint explicit. An applicant has an unobserved type $T\in\{g,b\}$ with
equal probabilities. Conditional on action, the potential success indicator
$Y(1)$ has probabilities
\begin{equation}
 P\{Y(1)=1\mid T=g\}=\frac45,\qquad
 P\{Y(1)=1\mid T=b\}=\frac15.
 \label{exec19:feast-law}
\end{equation}
The joint probabilities for $(g,1),(g,0),(b,1),(b,0)$ are respectively
$2/5,1/10,1/10,2/5$. At decision hour $10$, both alternative worlds have the
same observed customer key and the same available score $1/2$; the type label
used to describe the generating law is not an input to the decision. A
successful action has stipulated payoff USD~100, every action costs USD~50,
and no action has value zero. The action value in each type is consequently
\begin{equation}
 v_g=100\frac45-50=30,\qquad
 v_b=100\frac15-50=-30.
 \label{exec19:feast-value}
\end{equation}
These amounts are assumptions of the illustration, not observed bank prices
or estimates from Feast.~\citep{exec19_feast_reference}

The original score has event hour $6$, creation hour $7$ and availability hour
$7$. A correction reveals the type-specific score, retains event hour $6$,
and has creation and availability hour $11$. With a four-hour lifetime, the
original lies exactly on the admissible lower event-time boundary. An
event-only retrieval can select the correction despite its later creation;
creation filtering selects the original. In a second history, add a record
with event hour $9$ and creation hour $9$ that reveals the same type-specific
score but is not delivered to the decision system until hour $12$. Now both
event-only and creation-filtered retrieval select that unavailable record.
An actual-availability predicate still selects the original. No timestamp
flag can compensate for a field whose meaning is different from the required
clock.~\citep{exec19_feast_reference}

Let the proposed action rule take action precisely when its retrieved score
exceeds $1/2$. The exact comparison is
\begin{center}
\begin{tabular}{p{0.39\linewidth}rrr}
\toprule
History & Event only & Creation cutoff & Availability cutoff\\
\midrule
Correction created after decision & USD~15 & USD~0 & USD~0\\
Additional delivery delay & USD~15 & USD~15 & USD~0\\
\bottomrule
\end{tabular}
\end{center}
The USD~15 entries equal $P(T=g)v_g=(1/2)30$: they evaluate a rule that knows
the type and acts only for $g$. That rule is unavailable at decision hour
$10$. Every feasible rule in this example has the same action probability
$r$ in both worlds, because its available inputs and independent
randomization reveal no type. Hence
\begin{equation}
 E\{\text{net value of a feasible rule}\}
 =r\left(\frac12v_g+\frac12v_b\right)=0.
 \label{exec19:feast-feasible}
\end{equation}
The apparent improvement is thus an information advantage that cannot be
implemented at the declared clock. This derivation does not estimate an
upstream software error or prove a general monetary loss from delayed data.
It supplies a reproducible reason to include availability checks in a
public-data prototype before funding bank experiments.~\citep{exec19_feast_reference}

The reference evaluates the entire four-state law with rational arithmetic
and checks invariance to source-row permutations. Its boundary cases separate
the lifetime endpoint, an earlier expired record, a future event, all-late
candidates, unknown clocks, observed nulls, observed zero and conflicting
equal-timestamp values. In an actual cohort, the adapter must also preserve
one result per request identifier. Dropping an applicant because all features
arrive late changes the evaluated population; treating the missing value as
zero changes its meaning. Feature unavailability can itself depend on
customer behavior or servicing workload, so the analysis must address the
resulting selection mechanism explicitly.~\citep{exec19_feast_reference}

\paragraph{Historical retrieval, online access and agent tools.}
The selected online API retrieves the latest online feature values, and the
online-store documentation describes materialization into a latest-value
store. A successful online read is therefore not evidence that an old
decision's information has been reconstructed. There is also a separate
provenance obligation for transformations: a feature carrying a predecision
timestamp can still use future observations, a retrospectively fitted
normalizer or a later model version. The proposed library result must bind
the input snapshot, feature and transformation versions, provider, retrieval
mode and relevant clocks. A CLI or MCP wrapper should preserve that same
meaning.~\citep{exec19_feast_source}

Feast's selected source distinguishes two MCP surfaces. The standalone
\texttt{feast mcp} process proxies feature and registry servers. Its feature
tools include online retrieval and vector search, but also push and
materialization operations. The selected client turns HTTP error statuses
into tool errors, rather than presenting a denial body as feature data. The
in-process MCP endpoint uses a different extra and deployment path. These
facts justify an integration design, not an assertion that the proposal's
historical contract already exists in either surface.~\citep{exec19_feast_source}

For the proposed causal workflow, an agent should request a separately named
historical operation carrying decision identifiers, declared clocks, feature
versions and failure states. Registry browsing can help locate candidate
features; it cannot certify their causal role or availability. Online reads
serve current decisions, while push and materialization change system state.
Those capabilities need separate operational contracts. An unsupported
historical filter, HTTP failure, cancelled job or missing feature must remain
a typed failure or missing observation throughout the library, CLI and MCP
path. Retrying with weaker time semantics would silently answer a different
question.

\paragraph{Reuse choice, alternatives and a bounded public qualification.}
The proposed route is to retain the original Feast implementation behind a
typed historical boundary, initially with one selected provider. The root
source declares Apache-2.0, including redistribution and notice obligations.
The selected project requires Python~3.10 or later and declares thirty direct
dependencies with version ranges; provider extras and the two MCP extras add
different dependencies. The source version is dynamically determined, so a
commit identity alone should not be presented as a verified package-release
number. A complete distribution, resolved environment and applicable
dependency notices remain necessary before executable admission. No native
provider or MCP conformance is established here.~\citep{exec19_feast_source}

An explicit bitemporal file or database query is the reference alternative
for a small public prototype: it makes both clocks and the revision rule
visible with fewer components, although its scale and operating costs remain
to be measured. Existing bank feature infrastructure may eventually be the
better integration choice if it already retains the required history and
availability evidence; that cannot be inferred from a feature-store label.
Porting Feast would duplicate provider, timestamp, conversion and failure
semantics without a demonstrated platform requirement. The conditional
preference is therefore integration of an original provider after
qualification, with porting deferred.

Public or generated histories can exercise late corrections, ingestion
delays, null clocks, lifetime boundaries, tied revisions, duplicate requests
and provider errors. The test must compare complete rows and missing states
against a frozen reference, then check that library, CLI and MCP responses
preserve the same retrieval mode. The data owner is responsible for the
meaning of availability and revision fields; the scientific owner for the
information set, population and causal assumptions; and the adapter owner
for conversions, row conservation, failure handling and version changes.
Costs to measure include history retention, joins, validation, data transfer,
materialization and upgrade regressions. These public checks can establish
whether an expensive bank trial is worth proposing. They do not supply bank
truth, measured commercial prices or institutional acceptance.


% SOURCE docs/release/2026-10-29/source/docs/writing-repair/2026-10-29/src/6ac706a0b4c5289b.tex CHAR 0
\section[Observation data and computation workflows]{Observation data and computation workflows across five language ecosystems}
\label{sec:exec25-dww}

A bank selecting an observation-data engine or computation workflow must preserve the meaning of an observation before it optimizes its movement. A loan, a customer, a report about that customer, a proposed linkage and an execution attempt are different objects. Conflating them changes the population, information set or denominator of a causal analysis. This section compares selected implementations of data.table and targets in R, Dask and Snakemake in Python, Beam in Java, Arrow and Taskflow in C++, and DataFrames.jl and Dagger.jl in Julia. The separate detailed DuckDB, Polars, Feast, Pandera, Splink, OpenLineage, Ray and Spark comparisons remain applicable under their original limits. Some follow this overview in the assembled chapter; the comparison here does not replace their source-specific arguments. R interfaces to C++ engines and Python interfaces to shared engines must retain both identities; the interface language does not identify the entire implementation.~\citep{exec25_dww_datatable,exec25_dww_targets,exec25_dww_dask,exec25_dww_snakemake,exec25_dww_beam,exec25_dww_arrow,exec25_dww_taskflow,exec25_dww_dataframes,exec25_dww_dagger,exec19_duckdb_source,exec23_polars_sources,exec25_otm_spark}

The comparison is a proposal for integration and public qualification. The nine new repository selections supply 83 authentic source files, including seven implementation delegates, pinned at or before 7 October 2026. Five bounded discovery queries do not establish ecosystem saturation. Unreturned search results, uninspected candidates, other languages, full dependency graphs and installed artifacts remain outside the established coverage. No downloaded implementation was executed. The finite examples below are project constructions whose arithmetic reproduces; they are neither observed package incidents nor measured bank losses. Package documentation's scholarly references are leads unless their genuine originals have separately been admitted.~\citep{exec25_dww_choice}

\subsection{The observation is a versioned scientific object}

Write an observation as
\[
 o=(u,r,v,t_e,t_a,t_d,m,\pi,\ell),
\]
where $u$ is the proposed scientific unit, $r$ the source record, $v$ its version, $t_e$ event time, $t_a$ availability time, $t_d$ the decision time, $m$ the missingness or measurement state, $\pi$ a justified inclusion probability when available, and $\ell$ the linkage information. This is a proposed interface contract, not a claim that a data frame contains all these fields. A historical replay for decision $t$ must use a declared information set, for example records with $t_e\leq t$ and $t_a\leq t$, together with a specified rule for revisions. A source's event timestamp alone does not prove that its contents were knowable then. A schema validator can enforce field types; it cannot manufacture assignment, inclusion or linkage truth. The Feast, Pandera and Splink treatments in this chapter supply complementary feature, validation and uncertain-linkage contracts.~\citep{exec19_feast_source,exec19_feast_reference,exec24_pandera_source,exec24_pandera_reference,exec24_splink_sources,exec24_splink_contract,exec25_dww_choice}

Consider a credit observation with event time 2, availability time 3 and value 20, followed by a revision with event time 4, availability time 8 and value 200. At decision time 5 an event-only backward lookup selects 200, whereas the two-clock replay selects 20. Both can be perfectly sorted and computationally reproducible. Their 180-unit difference is an information error, not a sampling standard error. The revision is useful for later decisions and retrospective descriptive reporting, but it cannot silently enter a decision-time causal feature. A proposed adapter therefore retains event time, first availability, revision lineage, clock units and timezone conventions, and treats absence of availability history as an explicit limitation.~\citep{exec25_dww_choice}

Arrow's selected as-of implementation illustrates why the sign of a tolerance is part of the scientific request. For left time $t$ and tolerance $h>0$, the inspected predicate accepts right times $s$ with $t\leq s\leq t+h$. A nonpositive tolerance instead selects the backward direction within its magnitude. Thus a left time 5, right time 7 and tolerance 2 can be a valid forward match. Positive tolerance also changes memoization to retain future queues. The selected advance method refuses decreasing on-key times; this ordering check does not enforce availability time. We would adopt Arrow's typed columnar representation and consider its joins behind an adapter that rejects an unspecified clock/direction contract. We would not port its optimized join merely to obtain a different language label.~\citep{exec25_dww_arrow,exec25_dww_choice}

\subsection{Cardinality, missing identities and linkage uncertainty}

Let a join map records from tables $L$ and $R$ through key $k$. The output count is
\[
 N_{\rm join}=\sum_k n_L(k)n_R(k).
\]
It counts matched pairs, not necessarily persons, loans or independent observations. With two loan rows having the same key and exposures USD~100 and USD~200, and two contact rows sharing that key, an all-match join produces four rows and a naive exposure sum of USD~600 instead of USD~300. Four equals $n_L+n_R$, so a guard rejecting only counts strictly greater than that sum does not reject this example. The construction demonstrates a logical limit of an expansion guard, not an assertion that the package promised unique financial units.~\citep{exec25_dww_choice}

In the inspected data.table development source, the ordinary join signature defaults to all matches, retained unmatched rows and no rolling. The option defaults set the many-to-many compatibility option true and the cartesian expansion permission false. In the selected duplicate-match branch, the row limit is $nrow(x)+nrow(i)$; the C delegate rejects strictly above that limit and separately guards its integer-size limit. An explicit alternative many-to-many policy exists in the inspected source. An adapter must record these options and assert the scientific cardinality independently: one-to-one for a unique identity crosswalk, many-to-one for a declared lookup, or weighted pair contributions for a probabilistic linkage. Calling a table keyed or sorted is insufficient.~\citep{exec25_dww_datatable}

DataFrames.jl exposes a different useful contract: explicit left and right uniqueness validation invokes duplicate-key checks. Its default missing-key error can be replaced by equality or unequal treatment. The selected implementation filters missing keys differently for left, right and inner joins under unequal treatment and refuses that mode for an outer join. Numeric NaN and negative-zero keys are refused in the inspected branch. These choices cannot be hidden inside a universal null conversion. Two unknown identities on each side can create four pairs if all missing keys are treated as equal; none is thereby verified as a person match. A reason code for absent identity is not a shared identity.~\citep{exec25_dww_dataframes,exec25_dww_choice}

Suppose a proposed linkage assigns probability $p=4/5$ to a loss of USD~100 and probability $1/5$ to USD~1,000. Conditional on this illustrative linkage model,
\[
 E[L]=280,\qquad
 \operatorname{Var}(L)=p(1-p)(1000-100)^2=129600.
\]
A top-match plug-in of USD~100 loses USD~180 of expected loss and all of this conditional linkage variance. Across 10,000 hypothetical accounts the expected-loss understatement is USD~1.8 million; no independence assumption is needed for that expectation calculation, but portfolio variance would require the dependence law. Splink-style match information should therefore survive data extraction and reach downstream sensitivity or integration, rather than becoming an unexplained hard join. Public qualification should include known finite match assignments, uncertain matches, unmatched records and shared identifier collisions. The underlying match probabilities themselves still need calibration evidence.~\citep{exec24_splink_sources,exec24_splink_contract,exec25_dww_choice}

The resulting R decision is conditional adoption of data.table for its bounded tabular role through explicit cardinality, type and time checks. The Julia decision is analogous for DataFrames.jl, with explicit uniqueness and missing-value modes. DuckDB remains a separately audited SQL/embedded-engine alternative and Polars a separately audited expression/columnar alternative. A wholesale rewrite would add conversion and maintenance risks without solving identification or linkage calibration. A small portable reference for cardinality and clock checks is more useful than an unqualified port of a mature engine.~\citep{exec19_duckdb_source,exec19_duckdb_reference,exec23_polars_sources,exec23_polars_reference,exec25_otm_polars,exec25_dww_choice}

\subsection{Task identity, cache identity and scientific identity}

For a requested estimate, distinguish a scientific request identifier $q$, a logical task $j$, and its physical attempts $(j,a)$. A proposed cache identity includes the exact input snapshot and schema, causal target and assumptions, transformations, numerical configuration, source/artifact versions and relevant random-state policy. A job name or a checksum of only the function body is not the complete identity. If an external table changes behind an unchanged filename, a cache that does not bind the table snapshot can return a reproducible answer to an old question. This is a general integration risk to be tested, not a measured defect in every workflow engine.~\citep{exec25_dww_choice,exec24_openlineage_source,exec24_openlineage_reference}

The targets source makes the distinction concrete. Its thorough cue includes command, dependency, format, repository, iteration, file and seed comparisons. The builder first handles missing or error metadata; it then treats the always and never modes before later change comparisons. Consequently a selected never cue bypasses those later checks, but does not simply erase the earlier missing/error checks. Error modes have distinct continue, abridge, trim, stop, null and workspace branches; a null result must not be interpreted as a successful numerical estimate merely because it is an R value. The metadata schema includes warnings, errors, bytes and seconds. Those are valuable evidence fields, but they are not complete procurement, engineering or failure-recovery costs.~\citep{exec25_dww_targets}

Target-specific seeds depend on the configured seed and target name, while an inner command can override them and an NA setting disables automatic seeding. A name change can thus be scientifically relevant. More generally, if two attempts reproduce the same 100 independent unit-variance draws, concatenating their outputs gives 200 stored values but only 100 independent logical draws. The variance of their mean remains $1/100$; treating all stored values as independent gives $1/200$ and understates the standard error by a factor $\sqrt{2}$. Restart reproducibility and new Monte Carlo replication need distinct policies. We would integrate targets through a named R worker and preserve seeds, attempt identities, cue settings and error classes in the proposed result envelope.~\citep{exec25_dww_targets,exec25_dww_choice}

Dask offers a complementary Python task-graph role. The inspected tokenization code supports custom normalization and pickle/cloudpickle paths. Depending on the deterministic setting, some failures raise a tokenization error while fallback paths can return an object identity or random UUID. A generated task token therefore cannot be presumed a portable scientific provenance digest. The code also performs serialization round trips during tokenization, so an integration should accept a constrained typed request rather than arbitrary serialized objects from an untrusted caller. The inspected delayed-task guidance advises against input mutation. Its local task executor distinguishes failed payloads; that local path does not establish the behavior of every distributed scheduler, storage connector or retry plugin.~\citep{exec25_dww_dask}

Array reduction introduces another independent contract. Dask's selected sum forwards dtype, chunk summation and reduction-tree control. For the finite terms $10^{16},-10^{16},1$, the exact real sum is 1, while ordinary binary64 grouping can give
\[
 (10^{16}-10^{16})+1=1,\qquad
 10^{16}+(-10^{16}+1)=0.
\]
The authored scalar example demonstrates nonassociativity; it does not report a Dask benchmark. Deterministic scheduling, deterministic randomness and justified numerical error are separate properties. Public trials should fix the dtype and reduction policy, compare against exact or high-precision oracles, and record error as a function of cancellation, scaling and partitioning. Bitwise equality alone would reject some harmless differences and accept some consistently wrong computations.~\citep{exec25_dww_dask,exec25_dww_choice}

\subsection{Retrying computation and committing financial effects}

An acknowledgement can be lost after an external effect commits. If a task sends a USD~100 payment, loses the acknowledgement and repeats the payment, the logical instruction is one payment while physical executions can create USD~200 of effects. A durable idempotency key checked atomically with the effect, or an equivalent transactional protocol, is required for this particular guarantee. A local cache, a finished future or a retry counter does not independently prove such a protocol. Conversely, suppressing every retry after an ambiguous failure can lose legitimate work. The proposal should expose unknown outcome and reconciliation, not coerce ambiguity into success or failure.~\citep{exec25_dww_choice}

Snakemake's selected persistence implementation records code, input names, parameters, logs, software identity and stack hash when metadata retention is enabled. Its checksum collection drops unavailable checksums and is subject to a file-size policy. With metadata retention disabled, completion can remove incomplete markers and return without retaining that evidence. Its inspected scheduler handles an error, frees resources, increments the attempt and requeues while the restart budget remains; otherwise it marks failure. These mechanisms are useful for reproducible batch workflows but do not justify treating all output files or external effects as transactionally committed. We would adopt Snakemake for declared file-oriented research jobs with explicit environment, artifact and output-validation contracts. Dask remains the task/data-graph alternative; Ray's earlier actor/task review retains its own limits.~\citep{exec25_dww_snakemake,exec24_ray_sources,exec24_ray_contract}

The Java comparison must distinguish a dataflow model from a particular runner. Beam's selected Deduplicate transform is expressly bounded by time and window, compares encoded bytes and does not preserve order. Its defaults use processing time and ten minutes. The inspected stateful function emits an unseen value, sets a relative expiry, then clears its seen state when the timer fires. The same logical request can therefore require a separate persistent identity contract across windows and expiry. Beam's WithTimestamps documentation also distinguishes timestamps from windows: changing an element's timestamp does not automatically recompute its window, and permitted late timestamps can still be dropped downstream. Event time, processing time, availability time, watermark and financial settlement time cannot be collapsed into one field.~\citep{exec25_dww_beam}

We would consider Beam and the previously inspected Spark interfaces as conditional JVM dataflow workers, with a named runner, connector and state-store configuration. The public plan must test late/revised records, encoded identity, window changes, duplicate delivery and sink acknowledgements. It must report whether the sink protocol was actually qualified. We would reject the inference that a Java API or a successful DAG proves perpetual deduplication or correct causal population construction. OpenLineage is useful as a separate provenance interface; lineage events still require bindings to actual input and output artifacts.~\citep{exec25_otm_spark,exec24_openlineage_source,exec24_openlineage_reference,exec25_dww_choice}

Taskflow supplies a C++ task-graph alternative. Its inspected exception documentation routes blocking failures to the calling context and asynchronous failures through a returned future's get method. The executor tracks anchored exception state and retains the first caught exception in the selected path. Waiting for completion must therefore be distinguished from consuming the returned error, and one reported exception must not be labelled a census of every concurrently failing task. Cancellation also cannot reverse an already completed financial effect. We would adopt Taskflow only for a bounded worker role with explicit result/error consumption, per-attempt evidence and qualified external-effect behavior; we would not advertise distributed persistence from this local executor evidence.~\citep{exec25_dww_taskflow}

Dagger's selected worker-fault handler evicts results lost with a dead worker, clears cache references and invokes dependency rescheduling. The inspected delegate guards dependency registration, propagates an errored upstream, and increments running counts before exposing replacement work for scheduling. The source's own introductory caveat describes this as a last-ditch recovery path with residual robustness limits, and the handler requires surviving workers. These are reasons to qualify a particular pinned Julia deployment with worker loss, cached-result loss and failure propagation; they are not grounds for claiming a measured hang rate or rejecting every Dagger use. The comparison remains conditional because a static branch review cannot establish the complete concurrent state space.~\citep{exec25_dww_dagger}

\subsection{Language-specific decisions and commercial comparators}

The integration boundary should retain the mature engines and port only small contracts and finite oracles. R would combine a data.table or audited DuckDB data worker with targets when its cue and error policies fit the experiment. Python would compare audited DuckDB/Polars and validation/linkage/feature services with Dask, Snakemake or the previously reviewed Ray role. JVM deployments would compare Beam and Spark under named runners and connectors. C++ would use Arrow's representation or joins and Taskflow's local graph role without confusing them with a durable distributed workflow. Julia would combine explicit DataFrames.jl joins with a qualified Dagger execution profile; Scala and Rust retain the earlier Spark and Polars evidence, and other relevant ecosystems remain uninspected rather than absent. These are separate observation-data and computation-workflow decisions in each language.~\citep{exec25_dww_datatable,exec25_dww_targets,exec25_dww_dask,exec25_dww_snakemake,exec25_dww_beam,exec25_dww_arrow,exec25_dww_taskflow,exec25_dww_dataframes,exec25_dww_dagger,exec25_otm_spark,exec25_otm_polars,exec25_dww_choice}

The previously inspected historical SAS preview job-execution example is a useful bounded commercial comparator: it separates a job from its definition, polls state, distinguishes completed and failed, and exposes provider-specific results and separate logs or cleanup. The proposed open interfaces must likewise preserve job identity, terminal state, errors, artifact links and retention. That historical example does not establish a current supported server profile, entitlement, price, connector guarantees or an executed comparison. The historical MATLAB and HUGIN workflow-source omissions concern their particular requested automation contracts. The later manuals discussed in Section~\ref{gap29:software-followup} resolve other specific capability questions; they do not establish the missing workflow profile, and their cutoff qualifications remain explicit. Commercial comparison requires the exact edition, supported automation boundary, permissible deployment and full costs. A commercial workflow's existence neither makes the proposed scientific contracts redundant nor establishes our implementation parity.~\citep{exec24_workflow_contracts_sas,exec25_dww_choice}

The root licenses provide preliminary component distinctions: data.table declares MPL2 with its file alternative, Dask BSD3, Beam and Arrow Apache2, and the selected targets, Snakemake, Taskflow, DataFrames.jl and Dagger sources MIT-style terms. This is source-declaration evidence, not distribution clearance. Compilers, native libraries, plugins, runners, storage connectors and transitive packages require complete artifact inventories and appropriate rights review. Source dates and development metadata must not be presented as installed supported release versions. The proposed worker boundary limits coupling but does not by itself remove license obligations.~\citep{exec25_dww_datatable,exec25_dww_targets,exec25_dww_dask,exec25_dww_snakemake,exec25_dww_beam,exec25_dww_arrow,exec25_dww_taskflow,exec25_dww_dataframes,exec25_dww_dagger}

\subsection{A public qualification design before bank experiments}

The first public suite should enumerate scientific units and join pairs, include known uncertain links and missing identities, and generate event/availability revisions with an exact admissibility oracle. Each language receives the same logical fixture and an explicit serialization/type mapping. An adapter passes only for its declared cardinality, clock, null and linkage treatment; refusal on unsupported combinations is a testable result. Arrow's scanner batch and readahead settings illustrate why memory limits must cover buffered and in-flight data rather than only one batch. Performance comparisons must preserve the same scientific work and output precision.~\citep{exec25_dww_arrow,exec25_dww_choice}

The second suite should separate logical requests from attempts and inject failures before an effect, after the effect but before acknowledgement, during output materialization and during recovery. For an illustrative two-request run, each request has a failed attempt costing three cost units and a successful retry costing one. The completed-request denominator is two, total cost is eight, and cost per completed request is four. Reporting only successful-attempt cost gives one and discards three quarters of the incurred cost. Store requested, attempted, refused, failed, retried, completed and unresolved counts separately; a workflow status alone cannot supply this accounting. The units are authored accounting quantities, not measured CPU dollars or a procurement quote.~\citep{exec25_dww_choice}

A proposed CLI or MCP submission should bind a versioned request, immutable input references and checksums, the scientific contract, worker/artifact profile and a caller-supplied logical idempotency key. Submission returns an execution identity; status returns a documented state, attempt count and unresolved-effect indicator; retrieval returns typed outputs, validation and diagnostic records, with hashes for logs and artifacts. The design does not accept arbitrary executable snippets or serialized language objects as scientific requests. A repository adapter must declare its writable area, resource budget, cancellation behavior and result consumer. No sibling repository has been modified or integrated by this comparison, and these contracts are proposed rather than implemented protocol guarantees.~\citep{exec24_openlineage_source,exec24_openlineage_reference,exec25_dww_choice}

Scientific acceptance, numeric accuracy, runtime recovery, artifact custody and cost acceptance should have separate verdicts. A known route exists for the finite public fixtures, adapter schemas, complete dependency inspection and named runner/sink trials. There is no universal software repair for absent assignment truth, uncalibrated linkage probabilities, missing historical availability or arbitrary dependent-data uncertainty. Those require additional information or defensible restrictions. The available engines can move and schedule evidence efficiently; the proposed product must retain its scientific meaning through every transformation and attempt.~\citep{exec25_dww_choice}


% SOURCE docs/release/2026-10-29/source/docs/writing-repair/2026-10-29/src/a39850f68aade7fd.tex CHAR 0
\section{OpenLineage: execution history without automatic scientific acceptance}
\label{sec:exec24-openlineage-reuse}
A bank comparing collections policies needs to know which data, model,
assumptions and decision rule produced each result. That need makes a common
lineage vocabulary useful, but also creates a risk: a completed job can be
mistaken for a justified causal answer. OpenLineage supplies run, job and
dataset events with extensible facets. At the selected 7 October 2026 commit,
the inspected core schema identifies itself as version~2.0.2 and the Python
package declares version~1.54.0. These pinned sources support a concrete
reference choice for the proposal. They do not establish that an existing
server preserves every facet, that another workspace has adopted the
contract, or that any emitted claim is scientifically true.
\citep{exec24_openlineage_source}

\subsection{A valid event is not a complete analysis history}
The core schema requires event time, producer URI and schema URI. A run
event additionally requires run and job objects; the run requires a UUID,
and the job requires namespace and name. The description of event type
calls for a START and one terminal COMPLETE, ABORT or FAIL per run, but
event type is not in that object's required-field list. Input and output
datasets are optional too. These are distinct facts: descriptive protocol
requirements are not all enforced by single-document JSON validation.
\citep{exec24_openlineage_source}

Local authored fixtures validate the acquired declarative schema using an
installed JSON Schema validator with format checking and only internal
references. A run record without event type passes; a COMPLETE record without
inputs or outputs passes; an invalid UUID fails. Two identical COMPLETE
records each pass although the pair has no START. These results concern the
schema, not native client construction, server ingestion or stream validation.
The document validator cannot establish a unique, complete lifecycle by
examining one event at a time. A product adapter therefore needs a separate
history rule with attempt identity, permitted state transitions, duplicate
handling and reconciliation of missing terminal records.
\citep{exec24_openlineage_reference}

Dataset namespace and name identify a dataset, but the selected base schema
does not require a content hash or immutable snapshot. Two authored payloads
with the same namespace/name and different hashes show why this matters.
The proposed causal facet should bind the exact data snapshot, availability
clock, eligible population, event dictionary, target, graph or model,
identification assumptions, nuisance fits, split assignment and uncertainty
method. Large or restricted artifacts should be referenced by authorized
identifiers and digests, not copied indiscriminately into every event.
\citep{exec24_openlineage_source,exec24_openlineage_reference}

The base facet requires producer and schema URIs and permits additional
properties. An authored facet asserting that a result is identified passes
the selected base schema despite also declaring an unsupported claim. The
test does not dereference its custom URI: that URI is a field, not proof that
an additional schema or its scientific claims were checked. A production
consumer must validate an approved custom-facet schema and separately review
the evidence supporting the assertion. Serialization, schema conformance,
causal identification and bank acceptance are separate states.
\citep{exec24_openlineage_source,exec24_openlineage_reference}

\subsection{Delivery guarantees depend on the concrete client and transport}
The Java client documents a normal return after emission regardless of backend
exceptions. The inspected implementation delegates through a metric timer
without an exception guard. Its concrete HTTP transport performs a synchronous
POST and raises on response codes from 400 through 599; its response exception
is a client exception. The comment is therefore not a sufficient universal
no-throw contract. The circuit-breaker branch can also return without emitting.
Even an incremented emission-complete counter is not a scientific acceptance
receipt. The wrapper must distinguish suppression, attempted delivery,
acknowledged delivery, persisted lineage and accepted analysis.
\citep{exec24_openlineage_source}

The Java duration overload introduces another boundary. The abstract
transport's default implementation accepts a duration but simply calls the
ordinary emit method. The inspected HTTP transport does not override that
overload; it configures connection, socket and response timeouts separately,
defaulting to five seconds when no millisecond option is supplied. Thus a
caller's requested duration is not automatically an end-to-end delivery
deadline on this path. Other transports may provide stronger behavior; they
require their own source and runtime review. Do not describe every transport
as asynchronous or assume that a generic timeout parameter bounds all work.
\citep{exec24_openlineage_source}

The Python client adds environment, tags and source-location facets, can
reduce dataset identifiers before filtering, and can return early for absent
transport, disabled emission or a filtered event. Its HTTP transport posts,
checks response status and returns a response, while the public client emit
method returns no response object. The selected retry defaults allow HEAD
and POST and declare total/read/connect counts of five, a backoff factor of
0.3, and status codes 500, 502, 503 and 504. The transport passes this
configuration to urllib3 through Requests. Declared package bounds are not
a solved environment; the effective dependency behavior still needs a locked
profile. A retryable POST and a run UUID do not prove server-side idempotency
or durable exactly-once ingestion. \citep{exec24_openlineage_source}

These distinctions motivate a proposed adapter contract, not an assertion
that it has already been implemented. Preserve a durable local request and
artifact record; emit lineage with a stable logical run identity and explicit
attempt identity; and reconcile the receiver's acknowledgement with the
local record. A response lost after the receiver commits creates ambiguity:
blind repetition can duplicate work unless the receiver's deduplication rule
is known. Cancellation, transport timeout and scientific refusal need separate
fields. Sensitive custom facets also require a logging review because the
inspected clients contain debug event-serialization paths. No actual leakage
or deployment failure is inferred from that observation.
\citep{exec24_openlineage_source}

\subsection{A causal result facet and a request-complete denominator}
Suppose three authored analysis jobs finish their computation, but two
correctly abstain because their identification requirements are unmet. With
a stipulated total cost of USD~45, cost per completed job is USD~15,
whereas cost per accepted causal artifact is USD~45. Reporting only a job
COMPLETE count conceals that distinction. Correct abstention may be valuable;
it should remain explicit rather than disappear from cost or success
denominators. This is a designed accounting example, not a measured service
price or empirical adoption case. \citep{exec24_openlineage_reference}

The proposed facet should therefore state whether the output is an estimate,
identified set, sensitivity analysis, diagnostic, or structured refusal. It
should retain the target and horizon, assumptions and their evidence,
sampling/dependence unit, conditioning population, resolved estimator,
uncertainty object, numerical status, and exact reviewed artifact. An agent
can draft a causal question or graph, but its proposal is not promoted to
verified identification by being placed in an extensible JSON field.
The method owner, artifact reviewer and institutional approver retain their
separate responsibilities.

\subsection{Reference choice, integration alternatives and closure route}
The direct repository license is Apache~2.0. The Java build declares Java~8
source/target with a Java~17 toolchain and a concrete set of library versions;
the Python build declares Python~3.10 or later, bounded build tooling and
minimum runtime dependencies. These declarations support an environment
plan, not transitive rights clearance, a complete lock or a platform support
promise. The preferred conditional route is an original Java or Python client
with a reviewed causal facet and explicit acknowledgement protocol. Direct
schema-compatible emission is an alternative for another language, but it
transfers validation, transport and compatibility maintenance to this project.
A port of every integration is unjustified without a measured requirement.
\citep{exec24_openlineage_source}

For the present product proposal, retain OpenLineage as a versioned protocol
and client reference. Defer native admission until the method owner agrees
the causal vocabulary, the adapter owner freezes clients and receiver,
the rights owner reviews the complete distribution, and the evaluation owner
registers public/generated conformance trials. Trials should cover omitted
states, duplicated terminals, reordered events, suppressed emission, custom
facet rejection, receiver persistence, acknowledgement loss, retry/cancel
behavior and end-to-end costs. A local file lineage backend is an explicit
comparator for small prototypes. Integration with sibling repositories is a
future consumer-specific contract and test, not implied by their presence
on the machine. No sibling repository is modified by this proposal review.

This choice completes a bounded reuse argument while leaving delivery,
native adoption, causal truth and receiver-specific qualification open.


% SOURCE docs/release/2026-10-29/source/docs/writing-repair/2026-10-29/src/b0509f21613e49f3.tex CHAR 0
\section{OpenMetadata: catalog evidence and causal acceptance}
\label{sec:exec24-openmetadata}

OpenMetadata can organize evidence about a causal workflow without deciding
whether its causal conclusion is correct. A bank may already use a catalog to
locate loan snapshots, trace transformations, assign ownership and document
quality expectations. The product proposal should connect to that chosen
catalog when justified, rather than require a second institutional source of
ownership. The selected OpenMetadata commit is
\texttt{278e4abf6fb0c3027d5f1878bac92839995ef75c}, dated 7 October 2026.
Its platform declares 2.0.0-SNAPSHOT and its Python ingestion project declares
2.0.0.0.dev0 with Python at least 3.10. These identify inspected development
artifacts, not a qualified supported server release. The root license declares
Apache-2.0; the ingestion dependencies are separately maintained in setup
metadata. A usable bank distribution still needs a frozen environment, connector
scope and complete rights review. No server, client, connector or upstream test
was executed for this comparison.~\citep{exec24_openmetadata_source}

\subsection{Identity, lineage and the meaning of an edge}

The inspected add-lineage schema requires an edge with source and destination
entity references. Each reference requires an identifier and type; its name and
fully qualified name are optional. The reference is not an immutable data-snapshot
hash. An authored example binds one entity identifier to two different payload
hashes. That is not a catalog fault: the identity of a table and the identity of
its contents at an information cutoff are different. The proposed causal artifact
must retain both the catalog reference and the exact snapshot, extraction
query, availability cutoff and exclusion manifest.~\citep{exec24_openmetadata_source,exec24_openmetadata_reference}

Lineage details can carry a query, column relationships, pipeline, source class
and timestamps. Such relationships describe transformations and provenance.
For example, a collections score may depend computationally on delinquency,
previous contacts and a repayment table. This does not decide whether contacts
cause repayment, whether distress confounds contact assignment, or whether a
post-contact field was improperly used at time zero. A proposed agent can use
the catalog to retrieve candidate variables and their declared clocks, but
scientific edge proposals need a separate causal diagram and reviewed assumptions.
An LLM-generated description must retain its status as a proposal.
\citep{exec24_openmetadata_source}

Authored local schema fixtures accept a minimal edge with no lineage details.
They also accept unregistered causal-assertion fields inside the selected
lineage-details object, while rejecting an invalid UUID and an extra top-level
request field. Two opposite-direction edges each pass the individual request
schema. These findings concern declarative validation only: they do not prove
that a Java model retains unknown fields, that a server accepts a cycle, or that
any custom assertion is true. A single request schema is not a graph-acyclicity
or identification test. The complete basic schema also contains a Java-specific
regular-expression extension that fails Python's metaschema format check in the
local validator. The original bytes are retained; the successful fixtures are
restricted to exercised definitions with local reference resolution and UUID
format checking. Full cross-language schema compatibility remains unqualified.
\citep{exec24_openmetadata_source,exec24_openmetadata_reference}

\subsection{Authorization and read-back are part of the result}

The selected lineage PUT endpoint authorizes editing on both endpoint entities
and calls the repository before returning an empty HTTP 200 response. The
repository resolves nondeleted entity references and an optional pipeline,
processes timestamps, validates details and inserts a relationship. It then
updates search, cache, events and optionally RDF. Inspection of this sequence
does not establish a common transaction or complete delivery semantics across
all those stores. The proposed adapter must distinguish an acknowledged request,
a persisted relationship, a search-visible edge and a reconciled causal artifact.
\citep{exec24_openmetadata_source}

The inspected column validator filters invalid destination columns and removes
unrecognized source columns; it retains a column mapping only when valid source
columns remain. Consequently, successful completion does not by itself prove
that every submitted column relationship survived. In an authored abstract
trace, a request mapping sources balance and ghost to score retains only
balance; a second mapping to a missing destination is discarded. This trace
follows the selected filtering rule and is not a native server result. A consumer
requiring complete lineage should compare a read-back representation with the
submitted, versioned manifest and return a structured mismatch when necessary.
Silently treating partial retention as complete evidence would undermine the
scientific audit.~\citep{exec24_openmetadata_source,exec24_openmetadata_reference}

Catalog update time also differs from observation availability. The inspected
timestamp helpers choose among incoming, prior and server times; creation time
can retain the earlier value and update time the later value. These fields are
useful provenance, but do not establish when a business record became available
to a decision maker. The proposed request therefore binds the event clock,
observation clock, extraction clock and catalog update separately. A monthly
repayment table backfilled after a treatment cannot be made pre-treatment
information merely by associating it with an older catalog entity.
\citep{exec24_openmetadata_source}

\subsection{Retries, contracts and separate acceptance states}

The inspected Python client declares three retries, a base wait of thirty
seconds, retry status codes 503 and 504, and a 429 feature-limit branch. Its
default timeout is a connect/read pair of ten and three hundred seconds.
The retry loop uses increasing waits; an authored sequence of three retryable
responses followed by success incurs waits of thirty, sixty and ninety seconds,
or 180 seconds before request durations. This is an abstract configuration
calculation, not measured latency or a universal total deadline. Selected
transport exceptions are raised explicitly, whereas some unexpected exceptions
on a nonraw path are logged and can return null. An empty successful body also
returns null. The adapter therefore needs an explicit status and response
contract; a null value cannot serve simultaneously as failure and successful
persistence evidence. Retries require a separate duplicate and reconciliation
policy.~\citep{exec24_openmetadata_source,exec24_openmetadata_reference}

The data-contract schema includes a latest-result object with timestamp, status,
message and result identifier, and an extension field. These support a place
to describe quality outcomes; the schema alone does not establish how those
outcomes are computed or what a particular pass means. The selected OpenLineage
endpoint additionally checks enablement and authorization, invokes a mapper and
constructs an explicit response. This is a promising connection to the preceding
OpenLineage proposal. Its mapper, custom-facet retention, transaction behavior and
failure responses still require their own pinned qualification. Neither the
existence of a contract nor a successful lineage operation supplies a proof of
causal identification, uncertainty calibration or bank acceptance.
\citep{exec24_openmetadata_source}

The proposed causal extension keeps distinct states for metadata conformance,
snapshot identity, identification review, numerical qualification, artifact
acceptance and institutional approval. It binds the estimand, population,
treatment regime, graph and assumptions, evidence hashes, sampling law,
uncertainty procedure, refusal reasons and reviewer provenance. Generic catalog
status may summarize these fields for navigation but must not replace them.
For an agent translating a business question, a catalog match is a candidate
data source; it becomes eligible only after clocks, semantics and scientific
assumptions have been checked.

\subsection{A reference decision with a concrete prototype route}

The present decision is reference use. Integration remains disabled. The
preferred conditional route uses the bank's existing OpenMetadata service and
its original pinned ingestion/API interfaces, with a small typed causal-artifact
adapter and explicit read-back reconciliation. Direct REST emission is an
alternative when the worker language and deployment constraints justify owning
schema compatibility and transport behavior. An OpenLineage route is another
conditional alternative after its mapping preserves the required fields.
Immutable local evidence files remain a useful public-data prototype comparator.
A whole-catalog port or a mandatory second catalog is rejected for the current
proposal because the additional maintenance and migration cost lacks a
demonstrated requirement. These are proposal choices, not installed integrations.
\citep{exec24_openmetadata_source,exec24_openmetadata_reference}

Costs must retain their denominators. With an authored USD~60 total for four
submitted edges, three assumed persisted edges and one assumed accepted causal
artifact, costs are USD~15 per submission, USD~20 per persisted edge and USD~60
per accepted artifact. These stipulated counts and costs are not observed
OpenMetadata prices or bank returns. A real comparison includes schema and
connector maintenance, permissions, artifact storage, reconciliation, failed
requests, operator time and scientific review. The method owner specifies causal
acceptance; the catalog/adapter owner freezes the server and mapping; the rights
owner resolves the distribution; the evaluation owner preregisters generated
or public fixtures for malformed references, filtered columns, stale snapshots,
permission denial, retries, partial persistence and full costs. These steps are
known, while the future native and institutional outcomes remain unknown.
\citep{exec24_openmetadata_reference}

Source-level schema compatibility, complete dependency resolution, remaining
server delegates and chosen-consumer semantics remain explicit follow-ups.
They do not block independent proposal obligations. This reference decision
preserves a useful catalog option without promoting catalog lineage to causal
acceptance.


% SOURCE docs/release/2026-10-29/source/docs/writing-repair/2026-10-29/src/d5b82c74a1f7c357.tex CHAR 0
\section{Polars as a data worker: population, information time and execution meaning}
\label{sec:exec23-polars-worker}

A banking analysis can preserve every numerical formula and still change its
question when it joins customer records to events. A customer with three
contacts can acquire three times the weight of a customer with one contact;
an apparently historical measurement can have arrived after the decision;
and an overlapping time window can be mistaken for an independent
observation. A dataframe engine is valuable because it makes these
transformations explicit and executable. It cannot decide which
transformation represents the banking experiment.

Polars is therefore considered here as a conditional data-worker reference.
Twenty-one selected files authenticate to revision
\texttt{bf2152618c7123836a825f8566f4d98585b6ff5b}, dated 7 October 2026 before
the Asia/Shanghai cutoff. The Python project declares version 2.0.0,
Python at least 3.10 and the matching \texttt{polars-runtime-32} dependency;
the Rust workspace separately declares 0.55.1 and edition 2024. These are
source declarations, not authenticated distributed releases. Runtime-64,
compatibility, Arrow, pandas, database, cloud and GPU routes have distinct
dependencies. The selected root contains MIT terms, including notice
retention, but this is not clearance of every dependency or commercial
service. No downloaded implementation or test is executed in this review.
\citep{exec23_polars_sources}

\subsection{A join changes weights unless the population contract prevents it}

The selected Python lazy interface defaults to an inner join and
\texttt{validate="m:m"}; the latter performs no uniqueness check. It offers
one-to-one, one-to-many and many-to-one validation, and forwards validation,
null equality, coalescing and order choices to its underlying join call.
The guide distinguishes inner, outer, semi, anti, cross and as-of joins.
These are substantive alternatives: a semi join filters membership without
adding the right-hand columns, whereas an ordinary equi join can multiply
rows. They must be chosen from the observation design, not from the fact
that a table happens to fit in memory.~\citep{exec23_polars_sources}

Let customer $i$ have outcome $Y_i$ and $m_i$ matching event records. For
one group with $n$ original customers and $\sum_i m_i>0$, the unweighted
mean after an ordinary join is
\begin{equation}
 \overline Y_{\mathrm{joined}}
 =\frac{\sum_{i=1}^n m_iY_i}{\sum_{i=1}^n m_i}
 =\overline Y+\frac{\operatorname{Cov}_n(m,Y)}{\overline m},
 \label{eq:exec23-polars-fanout}
\end{equation}
where $\operatorname{Cov}_n$ uses denominator $n$ and customers with no
match have $m_i=0$. The identity follows by expanding
$n^{-1}\sum_i(m_i-\overline m)(Y_i-\overline Y)$. Equal multiplicities
preserve the mean; dependence between outcome and multiplicity generally
does not. Neither the identity nor a successful join establishes a causal
effect.~\citep{exec23_polars_reference}

For an exact supplied example, two treated customers have outcomes 10 and
0 and two controls have outcomes 0 and 0. Their customer-level descriptive
contrast is 5. If the first treated customer has three event records and
each other customer has one, the joined treated outcomes are
$(10,10,10,0)$ and the contrast becomes $15/2$. Against a stipulated
cost threshold of USD 6, the two descriptive calculations imply margins
of USD $-1$ and USD $3/2$. That reversal illustrates a population error;
it is not evidence that either contrast is a treatment effect or that
Polars has produced an incorrect join. An event-level target may be
appropriate, but it is a different target.~\citep{exec23_polars_reference}

The proposed worker should carry a sampling-unit key, an expected join
cardinality and pre/post membership counts. A many-to-one join requires a
justified unique right-hand key. If multiple rows are meaningful revisions,
aggregation or version selection must be scientifically specified before
joining. Dropping duplicates with an arbitrary first-row rule merely
hides the unresolved choice. Anti joins can expose unmatched IDs for
review, while a declared refusal can preserve an unresolved key rather
than silently removing the customer.

\subsection{Missing values and missing identities are different problems}

Polars distinguishes missing \texttt{null} values from floating-point
\texttt{NaN}. The selected guide states that \texttt{null\_count} excludes
NaN, \texttt{fill\_null} does not fill it, and numerical means and sums
skip nulls while NaN typically propagates. These are useful computational
contracts. They do not say that omitted outcomes are missing at random or
that a nonfinite number should be reclassified as an unobserved bank
response.~\citep{exec23_polars_sources}

Suppose the recorded outcomes are 2, 4 and one missing value. The observed
mean is 3. Filling the missing entry with zero changes the three-person
mean to 2. Two possible complete populations, with hidden values 0 and 8,
have the same observed data but full means 2 and $14/3$. Thus the
computational choice to skip, fill or reject does not identify the
full-population mean. A separate observation model or sensitivity range
must justify an estimand involving the missing response. These are exact
authored finite populations, not a claim about a package's imputation
accuracy.~\citep{exec23_polars_reference}

Join keys require another distinction. The selected interface defaults to
not matching null keys, while \texttt{nulls\_equal=True} permits them to
match. An inspected unit test describes two null keys on the left matching
a null key on the right; it has not been run here. Equality of missing
tokens is a declared relation on values, not evidence that two unidentified
customers are the same person. The worker must preserve unmatched and
ambiguous identities and their denominators. Treating every missing
identifier as one customer would impose a linkage claim that the engine
has no information to establish.~\citep{exec23_polars_sources}

\subsection{An as-of key must represent the information available to the decision}

The lazy as-of interface documents a backward search selecting the last
right-hand key at or before the left-hand key, a forward search selecting
the first key at or after it, and a nearest search. It defaults to
backward, permits exact matches unless told otherwise, and supports
grouping and a tolerance. Both inputs must be ordered within the relevant
groups. The same source qualifies checking: the in-memory engine cannot
check grouped sortedness, while the streaming engine checks the rows it
processes. Inspected tests describe these distinctions and deliberate
invalid sortedness declarations; they are intended checks, not observations
of this workspace's runtime.~\citep{exec23_polars_sources}

Let $t$ be the decision time, $e_j$ a measurement's event time and $a_j$
its availability time. A possible historical feature contract chooses
the largest $e_j$ among records satisfying both $e_j\leq t$ and
$a_j\leq t$, within the correct customer and permitted revision lineage.
Backward matching on event time alone imposes only the first condition.
At a supplied decision time 100, consider records
$(e,a,x)=(90,105,0),(80,85,4),(101,101,10)$. Backward event matching
selects value 0, although it arrived too late; nearest matching selects
the future value 10. The stated availability-aware contract selects value
4. The exact example demonstrates information leakage from the wrong
clock, not a defect in nearest or backward matching themselves.
\citep{exec23_polars_reference}

Equal timestamps need a declared ordering convention: a value logged at
the decision timestamp may have been recorded just after the action.
Disabling exact matches can express strict inequality but cannot recover
an unrecorded within-timestamp order. Tolerances likewise need units and
meaning. The selected source distinguishes calendar-day tolerances from
fixed 24-hour durations under daylight saving. Its time-zone guide
distinguishes converting a zone from replacing one. A wrapper should bind
the zone, precision, strictness, revision and availability clock instead
of allowing a serialization round trip to choose them implicitly.
\citep{exec23_polars_sources}

Window construction also enters uncertainty. The guide permits the
window period to differ from its start interval, allowing overlap or
gaps. For independent mean-zero, variance-one signs $X_1,X_2,X_3$, define
$W_1=(X_1+X_2)/2$ and $W_2=(X_2+X_3)/2$. Then
$\operatorname{Var}(W_1)=\operatorname{Var}(W_2)=1/2$ and
$\operatorname{Cov}(W_1,W_2)=1/4$, so
\begin{equation}
 \operatorname{Var}\!\left(\frac{W_1+W_2}{2}\right)=\frac{3}{8}.
 \label{eq:exec23-polars-overlap}
\end{equation}
Treating the two windows as independent gives $1/4$, understating variance
by a factor of $3/2$. Enumeration of the eight sign triples reproduces
these moments. A correct window aggregation therefore still requires a
dependence-aware uncertainty calculation for the resulting panel or time
series.~\citep{exec23_polars_sources,exec23_polars_reference}

\subsection{Streaming is a plan property, not a memory or recovery certificate}

The lazy interface exposes optimized and unoptimized plans. The guide
describes predicate, projection and slice pushdown, common-subplan
elimination, join ordering and type coercion. These can preserve the
declared relational query while reducing work; they do not determine
whether that query matches the experiment. The appropriate audit binds
the logical request and input identities as well as the selected physical
plan, engine and environment. An optimized graph is useful evidence of
execution choices, not evidence that a selected sample is causally valid.
\citep{exec23_polars_sources}

The selected \texttt{collect} documentation says that auto selection uses
configured engine affinity, otherwise streaming, and that unsupported
work can fall back to in-memory execution. GPU and background modes carry
explicit instability qualifications. The engine delegate forwards the
selected engine name to the underlying collect call; source comments
about relative performance are not measurements on this banking workload.
The streaming guide itself contains an editorial exclusion comment, so
its presence in the repository does not establish publication of that
page in the hosted manual.~\citep{exec23_polars_sources}

A documentation tension is more informative when followed into code.
The eager join docstring says validation is not supported by the
streaming engine; the lazy docstring lists validation and streaming
support without that same note. In the inspected Rust lowering path,
the specialized streaming join branch requires
\texttt{!args.validation.needs\_checks()}. Other cases produce an
\texttt{InMemoryJoin} node while retaining the join arguments.
\texttt{needs\_checks} is false only for many-to-many validation.
The inspected in-memory node waits for both input sinks, invokes its
join closure and classifies its sink state as accumulating memory.
This supports a precise source-level interpretation: requesting checked
join cardinality can introduce an in-memory join within a streaming
plan. It does not support saying that validation is silently discarded,
that all such queries fail, or that native execution has been verified.
\citep{exec23_polars_sources}

This tradeoff matters to a proposal's cost model. Removing a key check to
retain streaming changes the accepted-input contract. Moving that check
to a separate prepass adds I/O and requires the checked snapshot to be
the one subsequently joined. The future trial must measure complete
validation, scans, conversion, spilling, fallback, output and restart
costs against an identical accepted-output definition. Neither a runtime
flag nor a successful small fixture supplies a bank memory bound or a
durable recovery guarantee.

\subsection{A narrow reuse choice with substantive alternatives}

The present choice is reference, with a Python process worker as a
candidate prototype route and direct Rust use as a separate engineering
alternative. Reusing a pinned engine can preserve its rich join and
expression semantics. Porting that engine into R, Java or C++ would add
compiler, null, ordering and dependency obligations without demonstrated
benefit. Consumers in those languages can instead exchange explicitly
typed artifacts through the common CLI or MCP contract; that bridge is
proposed and still needs conformance evidence.

A relational database or another dataframe engine remains a substantive
alternative when its snapshot, key and query contracts are better suited
to the chosen workload. No speed ranking is assigned without a matched
trial. Small project-owned enumerations remain the independent oracle
for population, clock and dependence examples. The scientific owner
specifies units, population, assignment and information law; the adapter
maintainer binds schema, null/NaN, keys, ordering and runtime versions;
the operator measures full resource and recovery behavior.

Before native integration, freeze actual Python and Rust runtime
artifacts, optional engines, dependency notices, input snapshots, plans
and output validation. Register edge cases for duplicate and missing
keys, nonfinite responses, late revisions, ties, grouped sortedness,
overlapping windows and checked-join fallback. Preserve failures and
refusals and compare full accepted-output costs. Source review supports
the reuse decision; native execution, causal qualification and bank
adoption remain open.


% SOURCE docs/release/2026-10-29/source/docs/writing-repair/2026-10-29/src/fd021cf3ee170b81.tex CHAR 0
\section{Pandera: validation evidence, population changes, and the causal contract}
\label{sec:exec24-pandera}

A validation result is useful only when the bank knows what was actually checked.
Consider the monthly servicing table used to choose between an ordinary reminder
and a payment holiday. The table contains an episode key, decision time,
information-availability times, treatment, currency and subsequent repayment.
A dataframe validator can detect declared violations of this representation.
It cannot decide whether assignment was confounded, whether the analyst chose
the right population, or whether a seemingly historical variable was available
to the decision maker. Those questions motivate a narrower and more useful
role for Pandera: a candidate implementation of explicit observation contracts,
with its configuration and transformations included in the evidence.

The inspected source is the 7 October 2026 commit of the inherited
unionai-oss/pandera repository. Its project metadata points to the
pandera-dev/pandera identity, declares Python at least 3.10 and obtains the
version dynamically through setuptools-scm. The root license is MIT; the
selected pandas, Polars and CLI extras have their own dependency declarations.
The earlier PyPI observation of version 0.34.1 is retained as a separate
observation, not an assertion that a built wheel and this source tree are
identical. Dependency resolution, transitive rights and operational support
remain properties of a future concrete environment. The findings below come
from primary source inspection and project-authored examples, without running
the package or its tests.~\citep{exec24_pandera_source}

\subsection{A Boolean result conceals several different services}
Let $D$ denote the immutable input table, $S$ the declared schema, $e$ the
selected dataframe engine and $c$ the effective configuration. A useful
validation record is closer to
\begin{equation}
 (h(D),h(S),e,c,\mathcal C_{\rm run},\mathcal C_{\rm skipped},
   W,F,h(D'),\mathcal T)
\end{equation}
than to a single pass flag. Here $h$ is a content digest, $W$ records warnings,
$F$ failures, $D'$ the returned table and $\mathcal T$ the transformations
between input and output. A digest binds bytes; it does not certify their
scientific truth. This is a proposed product record, not a claim that Pandera
already emits every field.

In the inspected pandas entrypoint, disabling validation returns the supplied
object immediately. The scope decorator returns a passing check result when
configuration excludes that check's scope. The native pandas backend also
turns a failed check configured with \texttt{raise\_warning} into a warning
and a passing result. These are explicit modes with legitimate uses, but
absence of an exception does not distinguish them from completed mandatory
checks. A proposed causal worker must record the effective configuration and
reject disabled validation, skipped mandatory checks or warning-only mandatory
failures as evidence of full contract acceptance.~\citep{exec24_pandera_source}

Lazy error collection and lazy dataframe evaluation are different concepts.
The inspected error handler aggregates errors when its \texttt{lazy} argument
is true; otherwise it raises on collection. For the inspected native Polars
path, an unspecified validation depth defaults to schema-only on a LazyFrame
and to schema-and-data on a DataFrame. An explicit context or global depth
can override those defaults. The native Polars backend rejects declared
user-defined parsers, while the pandas backend runs its parsers. An optional
Narwhals flag changes backend registration. Consequently, a pandas review
cannot be silently transferred to Polars, a distributed engine or the
Narwhals implementation. The backend class, eager/lazy input and effective
depth belong in the request and result.~\citep{exec24_pandera_source}

\subsection{Sampling checks is not a census of the data}
The native pandas container applies parsers before its data-check subsample.
Its container metadata checks and sampled data checks follow different paths.
For head, tail or random sampling, the inspected helper concatenates the
selected objects and removes duplicate index labels. Whole-table claims
therefore require more than a successful sampled check.
\citep{exec24_pandera_source}

With $N$ records, $K$ violating records and a uniform sample of $n$ distinct
records without replacement, the probability of missing all violations is
\begin{equation}
 \Pr(\text{no sampled violation})=
 \frac{\binom{N-K}{n}}{\binom{N}{n}},\qquad n\leq N-K.
\end{equation}
For $N=100$, $K=2$ and $n=10$, the exact probability is $89/110$, approximately
80.9 percent. This calculation assumes unique observation identities and the
specified sampling law. It is not a confidence statement for an arbitrary
head sample, clustered records or the package's complete execution.
\citep{exec24_pandera_reference}

Index labels introduce a separate issue. In an abstract table with
$(\text{index},\text{amount})$ rows $(7,10)$, $(7,-1)$ and $(9,20)$,
keeping the first two rows and then the first occurrence of each index leaves
only $(7,10)$. A negative-amount rule would no longer see the second record.
Conversely, deleting every row with failure label seven removes both records,
including the positive one. These are independently reproduced label
operations, motivated by the inspected subsampling and row-removal code;
they are not observed native package failures. The causal contract should
require unique immutable row identifiers and separately validate the banking
episode key. A dataframe index is not automatically either identity.
\citep{exec24_pandera_source,exec24_pandera_reference}

\subsection{Cleaning can change the question being answered}
The inspected pandas path supports adding missing columns, filling defaults,
coercing types and filtering columns outside the schema. With row removal
enabled, lazy error collection is required; eligible failures are removed by
their index labels, while the inspected string-valued schema-failure branch
still raises. These transformations must be explicit scientific decisions.
For example, integer conversion maps both textual account identifiers
\texttt{0012} and \texttt{12} to twelve. Filling a missing exposure field with
zero asserts something different from recording that its value is unknown.
Filtering an undeclared availability-time column can remove evidence needed
to reject future information.~\citep{exec24_pandera_source,exec24_pandera_reference}

The statistical effect of row selection can be large even when every
retained field has a permitted type and range. Let treatment $A$ and an
unobserved risk bit $U$ be independent fair binary variables, with potential
repayment $Y(a)=U$ for both actions. The causal effect is zero. Suppose a
quality or eligibility rule retains a record only when $A=U$, perhaps through
a proxy whose relation to treatment has not been examined. Exactly half the
population remains, and
\begin{equation}
 \mathbb E[Y\mid A=1,R=1]-\mathbb E[Y\mid A=0,R=1]=1,
 \qquad R=\mathbf1\{A=U\}.
\end{equation}
The retained observational contrast is one despite the zero causal effect.
The problem is the selection rule, not a claimed validator defect. A product
must bind the selection indicator and the intended target population, preserve
excluded records with reason codes, and require a justified missingness or
selection argument before using the reduced table.~\citep{exec24_pandera_reference}

Null handling also requires a declared policy. The inspected native pandas
column-check path may remove null values before invoking a custom check when
\texttt{ignore\_na} is true. Dataframe postprocessing depends on the return
shape: the inspected Series-output branch permits all-null rows, whereas the
table-output branch combines results with a cellwise null mask. The general
method description is not sufficient to infer identical behavior for every
check shape. Nullability constraints, the custom-check result and the missingness
mechanism are distinct. Empty or excluded input is not positive evidence of
an observed value satisfying the substantive condition.
\citep{exec24_pandera_source}

\subsection{Chronology and identification need their own evidence}
For a baseline feature $X_j$, distinguish event time $t_j^{\rm event}$,
availability time $t_j^{\rm avail}$ and the action time $t_A$. The useful
condition is generally $t_j^{\rm avail}\leq t_A$, with any permitted latency
and revision convention explicitly stated. An event on day one, an action on
day two and a record first available on day three satisfy event chronology
but violate availability chronology. A generated fixture should expose that
case before any expensive bank trial. A passing check still depends on the
truth of the supplied timestamps.~\citep{exec24_pandera_reference}

Even a complete, correct representation cannot identify every causal effect.
Let $U$ be a fair binary variable. The structural worlds
\begin{equation}
 \mathcal M_1:\ A=U,\ Y=A,
 \qquad
 \mathcal M_2:\ A=U,\ Y=U
\end{equation}
produce the same observational pairs $(0,0)$ and $(1,1)$ with equal probability.
Their average intervention contrasts are one and zero respectively. Every
deterministic check of the same observed table has the same answer in both
worlds. Thus a schema-valid table cannot resolve the causal ambiguity.
The appropriate product response retains the competing models or refuses the
unidentified query until justified assumptions or new information are supplied.
\citep{exec24_pandera_reference}

\subsection{Original worker, existing CLI, and proposed MCP service}
The preferred prototype candidate is a narrow original Python worker for one
frozen dataframe profile. Its request identifies authorized input and schema
artifacts, their hashes, engine and dependency versions, mandatory checks,
availability rules and allowed transformations. Its response distinguishes
execution failure, disabled or partial checking, contract rejection, transformed
output, and a separate scientific disposition. It returns the exact validated
artifact or its immutable binding; the caller must not assume that the original
input is the returned transformed table.

The inspected CLI already loads YAML or JSON schemas, selects their declared
API, loads data and calls validation with lazy error collection. It rejects
selected backend mismatches and catches validation exceptions with exit code
one. Its call does not assign the returned validated object, and the success
report follows a non-raising call. Therefore it is a useful comparator for
validation reporting, but source inspection does not establish an artifact
export contract, serialization of every custom check, or scientific acceptance.
Its YAML loader uses safe loading; the delegated schema deserializers and
report-generation internals remain separately unqualified. We have not run
the CLI.~\citep{exec24_pandera_source}

A proposed MCP tool should invoke that bounded worker using artifact identifiers
and a preapproved schema, not accept arbitrary agent-authored Python checks as
automatically trusted code. An LLM may draft candidate currency, key and clock
rules and explain failures. The method owner must review their meaning against
positive, negative and deliberately ambiguous public or generated fixtures.
The service should retain the schema proposal, review, actual configuration,
failure evidence and all output identities. This is an integration design,
not a claim of an implemented MCP server or a proved protection boundary.

Direct adoption is conditional on those contracts. A wholesale port into R,
Java or C++ is presently unattractive because dtype, null, index, parser and
backend semantics would all need independent replication and maintenance.
An existing bank validator remains a reasonable alternative if it satisfies the
same observation contract. Small project-owned checks can provide transparent
prototype oracles, but should not pretend to replace general dataframe engines.
Reference-only reuse preserves these design lessons while implementation and
qualification remain open. The direct license does not remove the cost of
environment maintenance or schema ownership.~\citep{exec24_pandera_source}

Qualification should compare original-worker and existing-validator paths on
the same immutable fixtures: duplicated row labels and episode keys, missing
currency, future availability, identifier coercion, warning-only failures,
disabled or reduced-depth checks, transformed output, and engine-specific
parser rejection. Register expected acceptance or refusal before execution.
Record every attempt and correct refusal, not just successful process exits.
In a stipulated twelve-attempt workflow costing USD 48, eight process successes
cost USD 6 each; if only three yield scientifically accepted artifacts, the cost
is USD 16 per accepted artifact. Four correct refusals are reported separately.
These are accounting examples, not measured prices or package benchmarks.
\citep{exec24_pandera_reference}

The method owner specifies the causal target and admissible transformations;
the data owner supplies trustworthy keys and clocks; the adapter owner freezes
the execution profile and failure contract; the rights and evaluation owners
resolve dependencies and preregister the public qualification study. Revisit
implementation after those contracts, environment and semantic oracles are
reviewed. The present decision retains Pandera as a qualified source reference,
without granting native integration, cross-engine equivalence or causal acceptance.


% SOURCE docs/release/2026-10-29/source/docs/writing-repair/2026-10-29/src/2ef7cde39e9a7a54.tex CHAR 0
\section{Splink and the causal cost of uncertain entity identity}
\label{sec:exec24-splink}
A bank linking servicing records to later outcomes must distinguish a plausible record pair from a trustworthy longitudinal observation. Splink supplies a concrete candidate for probabilistic linkage, comparison diagnostics and graph-based entity construction. These operations can support a causal workflow, but a matching score does not identify a treatment effect, certify a customer identity or supply the joint uncertainty of an entire linkage map. The proposed use therefore preserves candidate pairs, their selection mechanism and the information clock before considering any hard entity identifier.~\citep{exec24_splink_sources,exec24_splink_contract}

\subsection{Selected artifact and the boundary of the evidence}
The inspected repository revision is dated 6 October 2026 at 09:43:41 UTC, before the client cutoff. Its project configuration declares Splink 5.0.0, Python at least 3.10 and below 4, and direct dependencies on DuckDB, PyArrow and SQLGlot. Separate optional groups cover plotting, graph operations and additional database engines. The selected license file contains MIT terms. PyPI metadata independently identifies a 5.0.0 distribution uploaded on 28 September 2026; equality of the version string does not establish that the distribution has every byte of the later inspected source revision. No package was installed or executed. A deployment decision still needs an exact distribution, resolved dependencies, notices, supported database profile and maintenance owner.~\citep{exec24_splink_sources}

The source review authenticates selected files against the complete cutoff tree. It follows prediction, parameter training, blocking, threshold clustering, constrained clustering and evaluation delegates. Historical README and package metadata remain discovery evidence rather than substitutes for these technical paths. The review does not qualify every comparison function, database dialect, cache, persistence mechanism, optional extension or advertised scaling claim. Original execution behind a narrow adapter is a candidate; a wholesale port would require independently maintaining comparison semantics, SQL generation and graph algorithms without a demonstrated benefit.~\citep{exec24_splink_sources,exec24_splink_contract}

\subsection{A posterior score is conditional on a model and population}
Write $M$ for the event that two records represent the same entity, $\gamma$ for their comparison pattern and $\lambda=P(M)$ for the prior over a declared pair population. Under the selected comparison model, and suppressing term-frequency adjustments, posterior odds take the form
\[
 O(M\mid\gamma)=\frac{\lambda}{1-\lambda}\prod_j\frac{m_j(\gamma_j)}{u_j(\gamma_j)},
 \qquad P(M\mid\gamma)=\frac{O(M\mid\gamma)}{1+O(M\mid\gamma)}.
\]
Here $m_j$ and $u_j$ describe comparison outcomes among matches and nonmatches. The factorization needs its conditional-dependence assumptions; it is not a consequence of storing fields in separate columns. The inspected prediction path adds base-two log weights to the prior log odds, converts the result to a probability with separate positive and negative branches for numerical stability, and filters by an inclusive weight threshold. The null comparison level has zero weight. Selected term-frequency adjustments alter the weight using observed frequencies and configured floors. These numerical conventions must survive conversion to another interface.~\citep{exec24_splink_sources}

An authored exact calculation shows why the prior population matters. A likelihood ratio of 100 gives posterior probability $100/199$, about 0.503, at prior 0.01, but $100/1099$, about 0.091, at prior 0.001. The same observed agreement can therefore cross a decision threshold when the pair population or prior changes. If two fields are exact duplicates of one underlying signal with likelihood ratio 10, the correct one-signal calculation at prior 0.01 is $10/109$. Treating the duplicated fields as conditionally independent instead gives $100/199$. More agreeing columns have not created more information. These are specified probability examples, not empirical calibration measurements for Splink.~\citep{exec24_splink_contract}

The inspected prior-estimation method counts pairs admitted by deterministic rules, divides by a supplied recall estimate and then divides by the total possible comparisons. Its documentation requests rules with very few false positives, and its code checks the recall range and some count consistency conditions. Recall remains an external assumption. With 40 admitted pairs among 10,000 possible pairs, supplied recall one half produces prior 0.008; supplied recall four fifths produces 0.005. The calculation cannot validate the recall or the rules' truth from the admitted count alone. Boundary warnings should not be promoted into evidence that an extreme prior is scientifically appropriate.~\citep{exec24_splink_sources,exec24_splink_contract}

\subsection{Training design, stopping rules and unobserved comparisons}
The pinned training guide recommends directly estimating the prior and nonmatch probabilities before training selected match probabilities. Its rationale is explicitly conditional: random pairs should be overwhelmingly nonmatching, while blocking enriches the match population for training. The source API samples records and evaluates eligible pairs in chunks; its optional per-level count stopping rule is described as a heuristic rather than a precision guarantee. Seeds and database support have their own constraints. Reused records induce shared-pair dependence, so the number of processed pairs alone is not an independent-sample size for a confidence statement. A public prototype must preserve sampled record identities, pair inclusion, seed, chunking, early stopping and effective comparison counts.~\citep{exec24_splink_sources,exec24_splink_contract}

Training blocking also changes which comparisons can be estimated. The selected session removes comparisons involving blocking columns and adjusts the session prior with factors associated with the blocking rule. The API defaults to fixed nonmatch probabilities, while the session prior can change; using a trained session prior to populate the main model is a separate option that defaults to false. Thus a general recommendation to estimate the population prior directly must not be misread as a claim that every internal EM prior is held fixed. Different training blocks can supply complementary estimates, but agreement among them does not prove independence, correct selection or population calibration.~\citep{exec24_splink_sources}

The inspected EM loop compares the largest absolute parameter change with a tolerance and also imposes an iteration limit. Its final logging statement uses convergence language after the loop, including the path that exhausts that limit. This is a static control-flow observation, not an observed package failure. A proposed adapter must report whether the criterion was actually met, retain the parameter history and distinguish iteration exhaustion from successful convergence. An authored sequence of maximum changes $0.10,0.05,0.04$ with tolerance 0.01 never meets the criterion, even though three iterations have completed. Small parameter change, when obtained, is still not proof of global optimum or posterior calibration.~\citep{exec24_splink_sources,exec24_splink_contract}

\subsection{Blocking recall and the causal observation mechanism}
Prediction uses selected blocking rules or a manually registered blocked-pair table. The latter has its own chunking restrictions. Chunked prediction processes the designated combinations and combines outputs; a DuckDB-specific path prunes source inputs to records present in selected blocked pairs. Such execution paths require matched semantic tests before cross-engine or chunking equivalence is claimed. A pair excluded before scoring cannot be recovered merely by lowering a later probability threshold. The evaluated-label path explicitly distinguishes positives missed by blocking from errors in the scoring model. That distinction should be mandatory in the product's result contract.~\citep{exec24_splink_sources}

Suppose a public fixture contains 100 true links, of which blocking admits 60. The scorer retains 54 true links and six false links. Its recall among admitted true links is 90 percent and its precision is 90 percent, but end-to-end recall is only 54 percent. A report that conditions on scored pairs hides 46 missed true links. Labels sampled only from accepted pairs cannot identify the omitted recall; the evaluation design must include a justified sample outside the block, a known public oracle or an explicitly bounded sensitivity analysis. Annotation effort and unresolved labels belong in the cost denominator.~\citep{exec24_splink_contract}

For a banking illustration, consider two treatment arms, each containing outcomes 100 and 0 USD. Their complete-data contrast is zero. If the linkage mechanism retains only the high outcome in the treated arm and only the low outcome in the control arm, the observed contrast becomes 100 USD despite perfect identity for the retained pairs. This is selection bias, not a false-merge error. The example is a finite construction and supplies no causal identification claim. It motivates separate diagnostics for treatment-dependent linkage availability, outcome-dependent blocking and use of post-treatment contact fields. A customer address updated after a servicing intervention cannot silently become baseline identity evidence.~\citep{exec24_splink_contract}

\subsection{Pairwise probabilities do not define a joint entity map}
The selected threshold-clustering method constructs connected components from retained edges. It does not require every pair within a component to exceed the threshold. With scores 0.96 for $A$--$B$ and $B$--$C$, but 0.10 for $A$--$C$, threshold 0.95 joins all three records in an authored graph calculation; all six edge orders give the same component. This deterministic behavior is useful for explaining transitivity, but it cannot make the three scores a coherent joint posterior over equivalence relations. If the first two numbers were probabilities under one common posterior, then transitivity and the union bound require
\[P(A\sim C)\geq P(A\sim B)+P(B\sim C)-1=0.92.\]
The score 0.10 violates that requirement. The example does not assert that a real dataset generates those scores; it shows why independent pair labels and a joint identity distribution are different objects.~\citep{exec24_splink_sources,exec24_splink_contract}

The reviewed alternative single-best-link routine constrains selected datasets to at most one member per cluster and exposes distinct tie policies. Such a source-uniqueness constraint can be valuable, but a local mutual-best rule is not automatically a maximum-total-weight assignment. In an authored two-by-two score matrix with rows $(0.90,0.80)$ and $(0.85,0.10)$, choosing the mutually highest pair first and then the remaining pair totals 1.00; the alternative cross assignment totals 1.65. This is an objective comparison, not an execution result or an assertion that total score is the correct banking loss. One-to-one assumptions themselves fail for legitimately repeated records unless their scope is explicitly declared.~\citep{exec24_splink_sources,exec24_splink_contract}

A causal consumer needs either a justified joint distribution over feasible linkage maps, an identified set with sensitivity bounds, or a refusal when identity information is inadequate. If a specified joint model assigns probabilities three quarters and one quarter to two maps yielding contrasts $20$ and $-20$ USD, the model-averaged contrast is 10 USD and the between-map variance is 300 USD squared. If each conditional analysis has variance 4 USD squared, the law of total variance gives 304 USD squared. Reporting only 4 discards the identity contribution. These weights are stipulated joint hypotheses; they cannot be manufactured by independently drawing the package's pairwise probabilities. Calibration and causal validity still require the joint model, observation law and downstream assumptions.~\citep{exec24_splink_contract}

\subsection{Reference disposition and a testable integration route}
The completed proposal choice is to retain Splink as a reference and conditional prototype candidate, without granting native integration. A deterministic exact-key baseline should establish the value of probabilistic matching. A pinned original worker is preferable to a wholesale port when it can preserve the selected database and comparison semantics. Alternative linkage or joint-identity models should be compared on the same known maps, missingness patterns, causal targets and full loss function. Better pair discrimination alone is not an adequate selection criterion.~\citep{exec24_splink_contract}

The proposed CLI/MCP request binds immutable input tables, source/record identifiers, allowed feature clocks, link type, candidate-generation rules, parameter and term-frequency snapshots, declared seeds, thresholds, tie behavior, engine identity and a resource budget. The response separates candidate pairs, excluded/unscored status, estimated probabilities, clustering rule, proposed entity maps, convergence diagnostics and uncertainty limitations. Failed, refused and unresolved requests stay visible. A result must not label a blocked-out pair as a calibrated nonmatch or an estimated component as certified identity. Arbitrary SQL and database credentials require their own scoped execution boundary; this source review has contacted no live bank or MCP endpoint.~\citep{exec24_splink_contract}

Before admission, a public/generated study should freeze planted entity maps, duplicate fields, correlated missingness, treatment-dependent linkage, conflicting source identifiers, bridge edges, tie permutations, blocked-out positives, iteration exhaustion and engine/chunking profiles. It should compare pair accuracy, entity consistency, causal estimand error, refusal and complete cost. The scientific owner controls the target and linkage law; the data owner controls clocks and source truth; the adapter owner preserves conversion and failure semantics; the platform and rights owners qualify the selected distribution and environment. Authentic bank validation follows a useful prototype and remains a separate obligation.~\citep{exec24_splink_contract}

An authored cost example stipulates 40 USD computation, 160 USD annotation, 200 USD engineering and 80 USD storage/review, totaling 480 USD. Six attempts for four logical requests, three of which produce accepted artifacts, cost 80 USD per attempt, 120 USD per request and 160 USD per accepted artifact. These are explanatory inputs, not prices or a Splink benchmark. Revisit the reference decision only after measured need, permissible artifact identities and predeclared semantic, uncertainty, failure and cost criteria support a bounded integration. A faster matching engine cannot supply missing identity truth or a universally calibrated causal analysis.~\citep{exec24_splink_contract}

The chapter's reuse choices therefore follow the information each component preserves. A catalogue locates and describes evidence; preparation and feature tools construct the declared population and history; validation records which conditions were actually checked; linkage retains uncertainty about the unit; and lineage or scheduling records how the work proceeded. An existing engine behind a narrow reviewed interface is the conditional preference where it serves that role. Explicit local files and finite oracles remain the small-case comparators, while native qualification determines whether the additional machinery earns its complete cost. No single catalogue status, successful job or matching score can substitute for the combined observation and causal argument.

